\newcommand{\DoNotLoadEpstopdf}{}
\documentclass[11pt,reqno]{amsart}
\usepackage[T1]{fontenc}
\usepackage{lmodern,amsmath,amssymb,amsthm,mathtools,microtype}
\usepackage[margin=1.1in]{geometry}
\usepackage[hidelinks]{hyperref}
\hypersetup{pdftitle={Collision records, arithmetic, and raw local inversion in a triangular Lorentz gas},pdfauthor={Qian Qi}}
\setlength{\emergencystretch}{2em}
\allowdisplaybreaks[2]
\numberwithin{equation}{section}
\newtheorem{theorem}{Theorem}[section]
\newtheorem{lemma}[theorem]{Lemma}
\newtheorem{proposition}[theorem]{Proposition}
\newtheorem{corollary}[theorem]{Corollary}
\theoremstyle{definition}\newtheorem{definition}[theorem]{Definition}
\theoremstyle{remark}\newtheorem{remark}[theorem]{Remark}
\newcommand{\R}{\mathbb R}\newcommand{\Z}{\mathbb Z}\newcommand{\T}{\mathbb T}
\newcommand{\dd}{\,\mathrm d}\newcommand{\one}{\mathbf1}
\newcommand{\supp}{\operatorname{supp}}\newcommand{\TV}{\operatorname{TV}}
\title[Collision records and raw local inversion]{Collision records, arithmetic, and raw local inversion in a triangular Lorentz gas}
\author{Qian Qi}
\date{October 5, 2026. A2-DYN, revision 8}
\subjclass[2020]{37D50, 37A50, 60F05, 42A38}
\keywords{Sinai billiards, collision records, local limits, inducing, Fourier inversion, geometric response}
\begin{document}
\raggedbottom
\begin{abstract}
We study the joint displacement, collision count and flight time of
actual first returns in a triangular periodic Lorentz gas, uniformly
for disk radii in $[0.45,0.47]$. A convex optical action constructs true
periodic orbits and yields quantitative joint phase separation. Its
one-step thickening gives physical $L^p$ coercivity for locally
H\"older phases. We retain the exact critical-edge coefficient, raw
jump subtraction, and local Fourier inversion statements.

A smooth collision-space killing estimate controls the time $N_n$ of
the $n$th actual return by $\Pr(N_n>L)\le A e^{an-cL}$. This gives
exponential moments of the entire record on a fixed complex
neighborhood and a linear collision-count truncation budget even on
exponentially growing Fourier bands. We construct the genuine induced
twists on a common Lebesgue scale, prove their exact physical renewal
identity, and establish holomorphy and strong parameter continuity.
The moving-domain coarea results, exact Kac normalizations and
conditioned logarithmic bound for the largest unfinished block are
preserved. Together these results advance the original raw mixed-density
local limit problem. Its anisotropic phase reconstruction, covariance
and all-branch residual estimates are stated explicitly as the remaining
hypotheses of the local inversion theorem, not inferred from the
Lebesgue realization or the clock bounds.
\end{abstract}
\maketitle

\paragraph{The return record and its analytic realization.}
The object of the paper is the original, unsmoothed joint record of
physical first returns. The quantitative advance in this revision is
Theorem~\ref{thm:cumulative-return-tail}: a single collision-space
avoidance estimate controls the entire cumulative record, rather than
only one return block. Proposition~\ref{prop:linear-count-budget}
converts it to a linear physical-count budget. Theorems
\ref{thm:common-renewal} and \ref{thm:common-scale-holomorphy} then
identify the actual induced operators and their fixed complex domain.
Their distinction from an anisotropic spectral estimate is essential
throughout. All periodic, coarea, raw-edge and clock arguments from the
preceding revision are retained below.

\section{The mechanical question and its interfaces}
Let $a=(1,0)$, $b=(1/2,\sqrt3/2)$, $\Lambda=\Z a+\Z b$, and
\[
 I=[9/20,47/100],\qquad Q_R=(\R^2/\Lambda)\setminus\overline B(0,R).
\]
The speed is one and every reflection is specular. Three indices will be kept distinct: a physical collision count $N$, an induced return count $n$, and the number of experimental bits. An induced record has values in $\Z^2\times\Z\times\R$, with counting measure in the first three coordinates and Lebesgue measure in the last. A four-dimensional Lebesgue density would be the wrong object.

The geometric inverse of the preceding article identifies the table from two fixed collision fields \cite{A2G}. It supplies neither independent flights nor an induced transfer-operator estimate. Conversely, the stationary law of the equilibrium billiard is not the unknown preparation law of that inverse experiment. The distinction permits geometric estimates to be used for equilibrium statistics without first estimating the entire preparation law.

The joint arithmetic result is Theorem~\ref{thm:joint-periodic}. It concerns
the complete periodic phase of an explicitly enlarged physical first-return
section. Lemma~\ref{lem:excursion-family} constructs its infinite period
family, and Theorem~\ref{thm:period-excess} proves the strict excess
inequalities. These results go beyond the zero-roof determinant by ruling
out every nonzero roof frequency, without a numerical irrationality claim.
The real-rank and compact-frequency consequences are stated separately
from measurable cohomology regularity and operator contraction.

The original physical record, winding geometry, critical-edge and
conditional inversion arguments remain in the article. Theorem
\ref{thm:local-edge-LLT} refines the inversion requirement by using local
edge errors instead of small total edge mass. Propositions
\ref{prop:exact-induced-means} and \ref{prop:induced-suspension} give exact
normalizations and stationary endpoint control for the enlarged section.
The full raw local limit is not inferred from the arithmetic theorem:
its spectral, residual and weighted-insertion requirements are listed in
the final section. All probability laws here refer to actual specular
orbits, not independent flight surrogates.

The relevant prior local-limit theory includes the Lorentz-process theorem of Sz\'asz--Varj\'u \cite{SV} and the suspension framework of Dolgopyat--N\'andori \cite{DN}. Definition 2.2 and Proposition 6.3 of the latter concern a mixing local limit tested against fixed compactly supported functions under minimality and group assumptions. They do not state absolute integrability of every raw finite-time characteristic function, or the particular geometrically uniform density theorem sought here. Demers--Zhang \cite{DZ} give a perturbation framework for Lorentz-gas maps; its spectral continuity is not automatically twice differentiability in a moving geometric parameter. We use the corrected version of the flow-mixing work of Baladi--Demers--Liverani \cite{BDL}. No novelty is claimed for Fourier inversion, the Morse lemma, or the suspension identity as general methods.


Theorem~\ref{thm:return-stability} compares the actual moving return
events in density $L^1$ and first moment, at every fixed return count.
Corollary~\ref{cor:uniform-stationary-endpoints} then supplies uniform
stationary endpoint control over the whole radius interval. The proof uses
finite-itinerary stability, coarea and exact Kac means, not a hypothetical
common spectral space. These finite-record moduli do not purport to supply
the quantitative long-time estimates in the local-limit interface.

Theorem~\ref{thm:quantitative-excess} strengthens strict increase to
uniform two-sided exponential bounds, obtained from the exact nonlinear
half action. Theorem~\ref{thm:log-period-separation} uses their upper and
lower bounds together: the former selects periods of logarithmic length,
and the latter supplies a polynomial discrepancy rather than mere
nonvanishing. Corollary~\ref{cor:approximate-phase-bound} states precisely
the continuous-phase consequence. It does not assume the measurable
regularity needed for an operator application.

At the level of the clock, Theorems~\ref{thm:uniform-return-fluid} and
\ref{thm:uniform-physical-clock} turn fixed-record continuity and exact
means into a uniform functional first-order law. This includes the maximum
unfinished return over the observation window at scale $t$, whereas the
retained stationary endpoint result is at scale $\sqrt t$ for finitely
many deterministic endpoints. The distinction of scales is essential.
Corollary~\ref{cor:physical-rate-error} gives a geometric error term for
the physical rate without assuming independence of the radius estimate
and the observed orbit. No originality is claimed for the general ergodic,
subadditive or renewal inversion arguments; the point is their verified
application to the actual moving return family.

\section{Physical records and an exact marked identity}
Write an outgoing section point as
\[
 q=R n_\alpha,\qquad v=\sqrt{1-p^2}\,n_\alpha+p\,t_\alpha,
 \qquad (\alpha,p)\in\T\times(-1,1).
\]
The collision map is $T_R$, its next physical flight is $\tau_R$, and
\begin{equation}\label{eq:flux}
 \dd\nu=\frac{\dd\alpha\dd p}{4\pi},\qquad
 \bar\tau_R=\frac{\sqrt3/2-\pi R^2}{2R},\qquad
 \dd\mu_R=\bar\tau_R^{-1}\dd\nu\dd s\quad(0\le s<\tau_R).
\end{equation}
These formulas follow by integrating the boundary flux $\cos\phi\dd s_{\partial Q}\dd\phi$: total section flux is $4\pi R$ and phase volume is $2\pi(\sqrt3/2-\pi R^2)$. Reflection preserves this flux. Thus $T_R$ preserves $\nu$.

Each next disk has a lattice label $d\ne0$. With $D=d-q$ and $L=D\cdot v$, its candidate entry time is
\begin{equation}\label{eq:entry}
 \tau_d=L-\sqrt{R^2-|D|^2+L^2}.
\end{equation}
The physical branch is the positive entry root with the smallest time, not an arbitrarily selected root. Its endpoint normal is $n'=(q+\tau_dv-d)/R$, and its outgoing velocity is $v'=v-2(v\cdot n')n'$. These expressions define the label, point, both velocities, and time on every regular branch. On a compact finite itinerary separated from tangencies and competing roots they depend analytically on $R$ and the initial coordinates, by the implicit function theorem. This is a finite-itinerary assertion, not a common norm for all iterates.

Distinct disks have distance at least $g_R=1-2R\ge3/50$. A straight outgoing flight cannot return to its initial convex disk. In particular there is no finite-time collision accumulation. This family has uniform finite horizon: a rational lattice direction has row spacing at most $\sqrt3/2$, so every line in that direction meets an open disk when $R>\sqrt3/4$; an irrational line is dense on the torus. Compactness then excludes arbitrarily long free segments at the smallest radius. Increasing the radius preserves this upper bound. Tangencies and their finite preimages have section measure zero.

Put $S_j\tau(y)=\sum_{i=0}^{j-1}\tau(T_R^iy)$. For $y$ interpreted as the first future outgoing collision and $u$ as its arrival time, the first $k$ future collisions have times $u+S_j\tau(y)$, $0\le j<k$. Their points, lattice labels, and incoming and outgoing velocities are obtained by iterating \eqref{eq:entry}. The final incomplete flight has duration $t-u-S_{k-1}\tau(y)$.

\begin{proposition}[Marked stationary record]\label{prop:marked}
For every bounded measurable functional $\Psi$ of this record and its terminal flight,
\begin{align}\label{eq:marked}
 \mathbb E_{\mu_R}[\Psi;N_t=k]
 =\frac1{\bar\tau_R}\int\dd\nu(y)\int_0^{\tau(T_R^{-1}y)}
 &\one_{\{u+S_{k-1}\tau(y)\le t<u+S_k\tau(y)\}}\nonumber\\[-2mm]
 &\hspace{-8mm}\cdot\Psi(\mathcal R_{k,t}(y,u))\dd u,\qquad k\ge1.
\end{align}
The initial phase point is $\Phi_R^{-u}y$, so the entire initial residual flight is retained as well. For $t=11/100$ there are at most two future collisions, and \eqref{eq:marked} with $k=1,2$, together with the zero-collision case, gives the full marked law without an independence assumption.
\end{proposition}
\begin{proof}
Under \eqref{eq:flux}, set $y=T_Rx$ and $u=\tau(x)-s$. Invariance of $\nu$ gives the displayed domain for $u$ with Jacobian one. The $k$th and $(k+1)$st future collision times are exactly the two endpoints in the indicator. This proves the identity for arbitrary bounded measurable $\Psi$, first for indicators and then by integration. The event of equality at an endpoint has zero suspension measure. Three collisions require at least two complete intercollision flights after the first event; their total duration is at least $6/50>11/100$.
\end{proof}


\section{Explicit periodic geometry and a genuine inducing set}
Two elementary cycles have zero winding. The normal nearest-neighbor orbit has two flights, total length
\begin{equation}\label{eq:two-three}
 L_2=2(1-2R),\qquad
 L_3=3(1-\sqrt3 R)
\end{equation}
is the length of the inner equilateral three-cycle on centers $0,a,b$. Its contact angles are $\pi/6,5\pi/6,3\pi/2$ and its outgoing $p$ coordinates are all $-1/2$. The reflection cosines are respectively $1$ and $\sqrt3/2$. The equilateral segments lie inside their elementary triangle and outside the three vertex disks; any other lattice center is at distance at least $\sqrt3/2$ from that triangle. Hence all other disks are strictly avoided.

\begin{lemma}[A winding four-cycle]\label{lem:winding}
For every $R\in I$ there is an analytic angle $\theta_R\in(0.92,0.967)$ satisfying
\begin{equation}\label{eq:theta}
 f_R(\theta)=2R\sin\theta-\sin\theta\cos\theta
                       +\frac{\sqrt3}{2}\cos(2\theta)=0.
\end{equation}
The itinerary with centers $0,a,a-b,2a-b,a$ and normals at angles
\[
 -\theta_R,\quad\pi+\theta_R,\quad\theta_R,\quad\pi-\theta_R
\]
is a physical four-collision periodic orbit modulo $\Lambda$, with winding $a$. All its incidence cosines exceed $1/2$ and its clearance from every third disk exceeds $1/200$. Rotation through $\pi/3$ gives another cycle with winding $b$ and the same length.
\end{lemma}
\begin{proof}
Put $c=\cos\theta$, $s=\sin\theta$. Its first four contact points and the translated fifth point are
\begin{align*}
 q_0&=(Rc,-Rs),&q_1&=(1-Rc,-Rs),\\
 q_2&=(1/2+Rc,-\sqrt3/2+Rs),&q_3&=(3/2-Rc,-\sqrt3/2+Rs),\\
 q_4&=q_0+a.
\end{align*}
Two segments are horizontal. Write $d=2Rc-1/2$, $h=\sqrt3/2-2Rs$. The other two are $(d,-h)$ and $(d,h)$. The scalar equation $d\sin2\theta+h\cos2\theta=0$ is precisely \eqref{eq:theta}. Since $d,h>0$, it makes the diagonal direction the reflection of the horizontal direction at the intervening normal. The remaining reflection equations follow by symmetry, with cosine $c$ at every contact. The total length is
\begin{equation}\label{eq:fourlength}
 L_4=2(1-2Rc)+2\sqrt{d^2+h^2}.
\end{equation}

Here are reproducible rational bounds for the inequalities and absence of occlusion. Divide $I$ into 32 closed intervals of width $1/1600$. On each, bisect the endpoint equations \eqref{eq:theta} between $9/10$ and $1$ to width at most $10^{-10}$. The interval derivative
\[
 2R\cos\theta-\cos2\theta-\sqrt3\sin2\theta
\]
is less than $-1/2$. The endpoint signs and positivity of $\partial_R f$ show that the two endpoint root brackets enclose the root for the entire parameter interval. The resulting angles lie in $(0.92,0.967)$; $d,h>0$ and $c>1/2$ on every interval.

For clarity, the arithmetic of this finite certificate uses no floating root solver. Square roots are enclosed by integer square roots at denominator $10^{22}$. For sine and cosine use the Taylor polynomials through degrees 25 and 24 at the rational midpoint, with errors at most $|x|^{26}/26!$ and $|x|^{25}/25!$, respectively; add the interval radius, since both functions are 1-Lipschitz. All interval operations are rational.

For each segment check centers $ia+jb$, $|i|,|j|\le4$, excluding its two endpoints. The distance from the rational midpoint segment is computed by clamping the rational orthogonal projection to $[0,1]$. Subtract the endpoint and center enclosure radii and the largest obstacle radius. Each lower bound exceeds $1/200$. Every contact point has norm less than 3; a center outside the index box has norm at least $5\sqrt3/2$, so it too has distance greater than the obstacle radius from every segment. The complete deterministic calculation is specified in \texttt{tools/certify\_winding.py}; it obtains lower bounds greater than $0.00925$ for clearance and $0.569$ for incidence, and a derivative upper bound below $-0.731$. The signs and strict margins prove the assertions on the full intervals, not only at sampled radii. The implicit function theorem gives analytic dependence and uniqueness in the printed root interval.
\end{proof}

Choose $\epsilon=1/200$. In $(\alpha,p)$ coordinates let $Y_R$ be the union of the open squares of half-side $\epsilon$ centered at
\begin{equation}\label{eq:Y}
 (0,0),\quad(\pi/6,-1/2),\quad
 (-\theta_R,\sin\theta_R),\quad
 (\pi/3-\theta_R,\sin\theta_R).
\end{equation}
Angles are taken modulo $2\pi$. The squares are disjoint and
$\nu(Y_R)=4\epsilon^2/\pi$. Let $r_R$ be the first return time to $Y_R$ and $F_R=T_R^{r_R}$ on its recurrent full-measure domain, with invariant probability $\nu_R^Y=\nu|_{Y_R}/\nu(Y_R)$. Define
\begin{equation}\label{eq:induced}
 \kappa_R=\sum_{j<r_R}\kappa_{\rm coll}\circ T_R^j,
 \qquad \tau_R^Y=\sum_{j<r_R}\tau_R\circ T_R^j.
\end{equation}
Poincar\'e recurrence and invariance give this induced probability; they do not give an exponential return tail, a Markov partition, or a differentiable family of spectral projectors.

\begin{proposition}[Full augmented lattice at zero roof frequency]\label{prop:lattice}
The four cycles above give fixed points of $F_R$ with data $(\kappa_1,\kappa_2,r,\text{return period})$ equal to the rows of
\begin{equation}\label{eq:matrix}
 V=\begin{pmatrix}0&0&2&1\\0&0&3&1\\1&0&4&1\\0&1&4&1\end{pmatrix},
                  \qquad |\det V|=1.
\end{equation}
Consequently a circle-valued coboundary identity
$e^{i(u\cdot\kappa+s r)}=e^{ic}g/(g\circ F_R)$ that can be evaluated at these fixed points forces $u=s=c=0$ modulo $2\pi$.
\end{proposition}
\begin{proof}
Only one contact of each cycle belongs to \eqref{eq:Y}. For the four-cycle the successive $p$ coordinates are $s,s,-s,-s$. The same is true after rotation. The two positive-$p$ unselected contacts are separated from the selected centers; the smallest relevant angular separation is $2\pi/3-2\theta_R>0.16$. All negative-$p$ contacts are separated in $p$ from the selected winding contacts, and also from $0,-1/2$. The remaining contacts of the two- and three-cycles are separated by their displayed angles. Thus the physical first return counts are $2,3,4,4$, not freely assigned labels. The four selected contacts are interior to their squares, so the same finite return branches exist in neighborhoods. Direct subtraction of the first two rows gives determinant of magnitude one. Evaluation first gives $2s=c=3s$, hence $s=c=0$; the last two rows then give $u_1=u_2=0$.
\end{proof}

The qualification about evaluation is essential. A merely measurable transfer function is not automatically defined at a periodic point. A spectral application requires its own regularity or Liv\v{s}ic argument. Moreover, adding a roof twist $b\tau_R^Y$ to this proposition leaves $bL_j$ in every periodic equation. The four rows do not eliminate it. Joint roof nonarithmeticity and returned temporal nonintegrability remain distinct requirements. As one smaller unconditional conclusion, \eqref{eq:two-three} excludes a continuous cohomology $\tau_R=c+g-g\circ T_R$: the two periodic averages are $1-2R$ and $1-\sqrt3R$, which differ by $(2-\sqrt3)R>0$.


\section{A physical period family and the complete periodic annihilator}
\label{sec:joint-arithmetic}
The four records in Proposition~\ref{prop:lattice} settle the integer
coordinates at zero roof frequency. We now add a family of actual periodic
trajectories that settles every roof frequency. A slightly larger section is
specified explicitly; the original section and every assertion about it
remain as stated in the preceding sections.

\subsection{The enlarged section}
Put
\begin{equation}\label{eq:enlarged-section}
 U=\{(\alpha,p):|\alpha-\pi/6|<3/50,\ |p|<3/20\},
 \qquad Y_R^*=Y_R\cup U,\qquad c_*:=\nu(Y_R^*)=\frac{91}{10000\pi}.
\end{equation}
The rectangle $U$ is disjoint from the four squares defining $Y_R$.
At the square with the same angular center the $p$ coordinate is near
$-1/2$; the other three squares are separated in angle or in $p$.
The area of $U$ is $4(3/50)(3/20)$, and the section density is
$1/(4\pi)$, proving the value of $c_*$. Write
$F_R^*=T_R^{r_R^*}$ for the first return to $Y_R^*$, and let
$\kappa_R^*$ and $\tau_R^*$ be the sums of the physical displacement
and flight length before that return. Its invariant probability is
$\nu_R^*=\nu|_{Y_R^*}/c_*$. An assertion about $F_R$ does not become
an assertion about $F_R^*$ merely by changing notation.

Rotate coordinates through $-\pi/6$ and put
\begin{gather*}
 D=\sqrt3,\qquad O=(0,0),\quad C=(D,0),\quad
 A=(D/2,-1/2),\quad B=(D/2,1/2),\\
 d=\frac12-R,\qquad g=D-2R.
\end{gather*}
The rotated lattice consists of $(Dk/2,l/2)$ with integers $k,l$ of
the same parity. The normal orbit between $O$ and $C$ is unobstructed:
the two intermediate centers $A,B$ have distance $1/2$ from its
supporting line, and all other centers are farther away. It has two
physical flights of total length $2g$, zero winding, and exactly one
visit to $Y_R^*$ per period. Its selected state is $(\pi/6,0)$ in the
unrotated coordinates. The four cycles of Proposition~\ref{prop:lattice}
still have exactly one visit per period. Indeed, the equilateral cycle
has $p=-1/2$, the two-cycle has angles $0,\pi$, and the winding cycles
have $|p|=\sin\theta_R>0.79$. None gains a visit to $U$.

\subsection{A convex optical action}
For $-d\le y,z\le0$ define
\begin{align}
 x_R(y)&=\sqrt{R^2-y^2},\qquad u_R(y)=R-x_R(y),\nonumber\\
 e_R(y)&=\sqrt{(D/2-x_R(y))^2+(d+y)^2},\label{eq:excursion-lengths}\\
 \ell_R(y,z)&=\sqrt{(D-x_R(y)-x_R(z))^2+(z-y)^2}.\nonumber
\end{align}
The function $e_R$ is the length from the top contact
$q_A=(D/2,-d)$ of the disk at $A$ to a contact at height $y$ on the
facing arc of the disk at $O$ or $C$. The function $\ell_R$ is a
flight length between the two facing arcs. For $m\ge1$ let
\begin{equation}\label{eq:excursion-action}
 \mathcal A_{m,R}(y_1,\ldots,y_{2m})
 =e_R(y_1)+\sum_{j=1}^{2m-1}\ell_R(y_j,y_{j+1})+e_R(y_{2m}),
 \qquad L_m(R)=\min_{[-d,0]^{2m}}\mathcal A_{m,R}.
\end{equation}

\begin{lemma}[A uniformly regular physical family]\label{lem:excursion-family}
For each $m\ge1$ and $R\in I$, the minimum in
\eqref{eq:excursion-action} is unique. Its coordinates satisfy
\begin{equation}\label{eq:excursion-height}
 -\frac{2d}{5}<y_j<0,\qquad y_j=y_{2m+1-j}.
\end{equation}
The contacts with these heights form a genuine specular periodic orbit
with center word $A,(O,C)^m,A$. Every flight avoids all third disks.
The orbit has zero winding, $2m+1$ physical collisions, and exactly
$m$ returns to $Y_R^*$. The period length $L_m(R)$ is analytic in $R$
for each fixed $m$.
\end{lemma}
\begin{proof}
We supply the variational, mechanical and return checks separately.
On the whole box, $x_R(y)>11/25$, $x_R(y)\le47/100$, and
$1.73<D<1.74$. The positive horizontal quantities
\[
 a(y)=D/2-x_R(y),\qquad h(y,z)=D-x_R(y)-x_R(z)
\]
satisfy $a>79/200$ and $h>79/100$. The second derivative of
$-x_R(y)$ is $R^2/x_R(y)^3>0$. The Hessian formula for a Euclidean
norm shows that the Hessian of $e_R$ or $\ell_R$ is the sum of a
positive semidefinite term and the Hessian of its horizontal
component multiplied by its positive horizontal direction cosine.
The latter cosine exceeds $9/10$ on this box. Thus each segment
length is strictly convex in its variables; in fact the Hessian of
$\mathcal A_{m,R}$ dominates the identity. One may check the last
numerical bound from
$(9/10)(9/20)^2/(47/100)^3>1$. Compactness and strict convexity give
a unique minimizer.

No coordinate equals $-d$. For such a coordinate $y$ and an incident
interior segment, the derivative is
\[
 \partial_y\ell_R(y,z)=\frac{h(y,z)y/x_R(y)+y-z}{\ell_R(y,z)}<0,
\]
since $z\ge y=-d$. For an end segment,
$e_R'(-d)=-a(-d)d/(x_R(-d)e_R(-d))<0$.
Increasing that coordinate therefore decreases the action.
At an end coordinate equal to zero, $e_R'(0)=d/e_R(0)>0$, and
all other incident derivatives are nonnegative. At an interior zero
adjacent to a negative coordinate, its derivative is positive.
If a block of zero coordinates had no such neighbor it would reach
an endpoint. A variation into the box excludes all these cases.
Every coordinate is consequently in $(-d,0)$. Reflection in the
line $x=D/2$ reverses the center word and preserves the action, so
uniqueness gives the symmetry in \eqref{eq:excursion-height}.

Here is a uniform improvement of the height bound. A smallest
coordinate $y=-w$ cannot be interior, because both incident
derivatives there are negative. Take it to be $y_1$, using reversal
if needed, and write $z=y_2$. Stationarity gives
\[
 d-w=\left(a+e_R(y)\frac{h}{\ell_R(y,z)}\right)\frac{w}{x_R(y)}
           +\frac{e_R(y)}{\ell_R(y,z)}(w+z).
\]
The last term is nonnegative since $z\ge-w$. Moreover,
\[
 \frac{a+e_R(y)h/\ell_R(y,z)}{x_R(y)}
 >\frac{(19/10)(79/200)}{47/100}=\frac{1501}{940}>\frac32.
\]
It follows that $w<2d/5$. This bounds every coordinate.

Interior stationarity of the optical length is precisely equality
of tangential incoming and outgoing velocities at each contact.
The outgoing directions point into the exterior normal half-plane.
For example, at a contact on $O$, its outward normal is
$(x_R(y),y)/R$. Its dot product with a direction to $C$ has numerator
at least $x_R(y)h-|y||z-y|>0$; for a direction to $q_A$ both the
horizontal contribution and $y(-d-y)$ are positive. The checks at
$C$ are reflections of these. The symmetry of the two end contacts
implies the reflection law at $q_A$, whose normal is vertical. There
$e_R(y_1)<11/25$ and $d+y_1>3d/5\ge9/500$, so the normal
incidence cosine exceeds $9/220>1/25$. The other contacts are
uniformly transverse by the preceding horizontal-component bounds. Thus the stationary polygon
is specular rather than the straight-through solution of the
stationarity equation.

The long flights lie in $-2d/5<y<0$, so their distance from the disk
at $A$ is at least $3d/5\ge9/500$, and from the disk at $B$ at
least $d\ge3/100$. End flights join $q_A$ to the appropriate facing
arc and leave both endpoint disks in their exterior supporting
half-planes. Their third disk $O$ or $C$ has horizontal separation
at least $D/2>R$, and $B$ is vertically separated. All segments
lie in the rectangle
\[
       11/25<x<D-11/25,\qquad -d\le y\le0.
\]
For the remaining lattice centers, the possibilities $k=0,2$ with
nonzero $l$ have vertical separation at least $1-d\ge0.95$;
$k=1$, $|l|\ge3$ have separation at least $3/2-d$; and
$k\notin\{0,1,2\}$ have horizontal separation exceeding $1.3$.
This enumerates the whole lattice and proves nonocclusion, not just
a check of a finite list of nearby obstacles.

At an outgoing contact on $O$, the normal angle differs from
$\pi/6$ by at most $|y|/x_R(y)<1/22<3/50$. The outgoing velocity
points toward $C$ and its angular deviation from the horizontal is
at most $|z-y|/h<2/79$. Thus
\[
 |p|\le \frac1{22}+\frac2{79}<\frac3{20}.
\]
All $m$ contacts on $O$ belong to $U$. Contacts on $C$ have angles
near $7\pi/6$, and $q_A$ has angle $2\pi/3$; none belongs to
$Y_R^*$. The return counts are $2$ between successive contacts on
$O$, with one count $3$ through $C,A,O$. Their sum is $2m+1$,
and there are $m$ induced returns. The orbit is closed in the plane,
so its winding is zero. All selected contacts and unselected contacts
have strict membership margins, giving actual regular return branches.
Finally the positive Hessian and the interior minimum allow the
analytic implicit function theorem. This proves analyticity for each
fixed $m$; no bound for high parameter derivatives uniform in $m$ is
being inferred.
\end{proof}

\subsection{Strictly increasing bounded excess lengths}
The following argument supplies an infinite sequence of period relations.
It does not infer nonarithmeticity from finitely many decimal lengths.

\begin{theorem}[Strict period excess]\label{thm:period-excess}
For the family of Lemma~\ref{lem:excursion-family}, define
$E_m(R)=L_m(R)-2m g$. Then, for every $R\in I$,
\begin{equation}\label{eq:period-excess}
 0<E_m(R)<E_{m+1}(R)\le \frac{2d^2}{g}<\frac1{100}.
\end{equation}
In particular $E_{m+1}(R)-E_m(R)>0$ and these differences tend to
zero. The convergence of $E_m$ is used pointwise in $R$; the subsequent
compact-frequency assertion supplies its own uniformity argument.
\end{theorem}
\begin{proof}
Suppress $R$ and set $K(y,z)=\ell(y,z)-g$ on $[-d,0]^2$.
The horizontal component of the flight gives
\begin{equation}\label{eq:kernel-coercivity}
 K(y,z)\ge u(y)+u(z),
\end{equation}
with equality exactly when $y=z$. Define continuous value functions
\[
 V_0(y)=u(y),\qquad
 V_{j+1}(y)=\min_{z\in[-d,0]}\{K(y,z)+V_j(z)\}.
\]
They are nonnegative and vanish at zero. For $j\ge1$, choosing
$z=0$ gives $V_j(y)\le K(y,0)=O(y^2)$ near zero; $V_0$ has the
same bound. The implied bound is independent of $j$.
Equation~\eqref{eq:kernel-coercivity} implies $V_1(y)>V_0(y)$ for
$y<0$: equality at a minimizing $z$ would require both $z=0$ and
equality in \eqref{eq:kernel-coercivity}, hence $y=0$.

For any $y<0$, zero is not a minimizer of $K(y,z)+V_j(z)$. Indeed,
\[
 \partial_zK(y,0)=\frac{-y}{\ell(y,0)}>0,
\]
whereas $0\le V_j(z)=O(z^2)$. A sufficiently small negative $z$
therefore reduces the value. If $V_j>V_{j-1}$ on $[-d,0)$, evaluate
at a minimizer $z_*<0$ for $V_{j+1}(y)$. It follows that
\[
 V_{j+1}(y)=K(y,z_*)+V_j(z_*)
   >K(y,z_*)+V_{j-1}(z_*)\ge V_j(y).
\]
Induction proves the strict inequality for every $j$ and $y<0$.

The minimizing path is symmetric, and its middle edge has length
$\ell(y_m,y_m)=g+2u(y_m)$. Splitting the action at that edge gives
the exact identity
\begin{equation}\label{eq:bellman-excess}
 E_m=2\min_{y\in[-d,0]}\{e(y)-g/2+V_{m-1}(y)\}.
\end{equation}
The outer minimizer cannot be zero, since $e'(0)=d/e(0)>0$ and
$V_{m-1}(y)=O(y^2)$. Evaluating the formula for $E_{m+1}$ at its
negative minimizer, and using strict increase of $V_j$ there, proves
$E_{m+1}>E_m$. Since $e(y)\ge D/2-x_R(y)\ge g/2$, equality
$E_m=0$ would force $y=0$ and then $d=0$, which is excluded.
Taking all heights zero gives
\[
 E_m\le2\sqrt{(g/2)^2+d^2}-g
      =\frac{2d^2}{\sqrt{(g/2)^2+d^2}+g/2}\le\frac{2d^2}{g}.
\]
Here $d\le1/20$ and $g>79/100$, proving the last bound in
\eqref{eq:period-excess}. A bounded increasing sequence has summable
positive increments, which must tend to zero.
\end{proof}

\subsection{Joint arithmetic and its exact scope}
A regular periodic return orbit has a total record $(K,N,T,h)$:
winding $K\in\Z^2$, physical collision count $N$, flight length $T$,
and induced period $h$. Define its phase at
$(u,s,\beta,c)\in\T^2\times\T\times\R\times\T$ to be
$u\cdot K+sN+\beta T-ch$, with $\T=\R/(2\pi\Z)$.

\begin{theorem}[Trivial joint periodic annihilator]\label{thm:joint-periodic}
For every $R\in I$, if
\begin{equation}\label{eq:joint-periodic-phase}
 \exp i(u\cdot K+sN+\beta T-ch)=1
\end{equation}
for every regular periodic orbit of $F_R^*$, then
$\beta=0$ and $u=s=c=0$ on their tori. It suffices to test the four
cycles of Proposition~\ref{prop:lattice}, the long normal two-cycle,
and the family of Lemma~\ref{lem:excursion-family}.
\end{theorem}
\begin{proof}
The long normal cycle gives $\exp i(2s+2\beta g-c)=1$.
The $m$th excursion gives
$\exp i((2m+1)s+\beta L_m-mc)=1$. Dividing the latter by the
$m$th power of the former yields
\[
                    \exp i(s+\beta E_m)=1.
\]
Taking consecutive quotients shows that
$\beta(E_{m+1}-E_m)\in2\pi\Z$ for all $m$. If $\beta\ne0$,
Theorem~\ref{thm:period-excess} makes these quantities nonzero and
strictly smaller than $2\pi$ in absolute value for large $m$, a
contradiction. Hence $\beta=0$. The four unchanged first-return
records in \eqref{eq:matrix} now give $u=s=c=0$.
\end{proof}

\begin{corollary}[Finite uniform separation on compact frequency sets]
\label{cor:compact-periodic}
For every $B<\infty$ and $\varepsilon>0$, there are a finite list of
these physical periodic words and $\eta>0$ such that, uniformly in
$R\in I$ and in $|\beta|\le B$ at distance at least $\varepsilon$
from the trivial phase, at least one word has
$|\exp i(u\cdot K+sN+\beta T-ch)-1|\ge\eta$.
\end{corollary}
\begin{proof}
Each word has fixed integer record and continuous length on $I$.
Theorem~\ref{thm:joint-periodic} gives a nonzero discrepancy at every
point of the compact parameter--frequency set in question. The open
sets on which a given word has positive discrepancy cover that set.
Take a finite subcover and minimize the maximum of its discrepancies.
This minimum is positive and supplies $\eta$.
\end{proof}

\begin{corollary}[Full real period rank]\label{cor:real-period-rank}
The real vectors $(K_1,K_2,N,h,T)$ of five fixed cycles span $\R^5$.
The absolute determinant is $2(\sqrt3-1)$, independently of $R$.
Consequently no nonzero real linear combination of the four induced
record coordinates is cohomologous to a constant by a transfer
function for which periodic evaluation is valid.
\end{corollary}
\begin{proof}
Append lengths $L_2,L_3,L_4,L_4$ to the four rows in
\eqref{eq:matrix}, and append $(0,0,2,1,2g)$ as the fifth row.
Subtracting the first row from the fifth leaves only the last entry
$2g-L_2=2(\sqrt3-1)$. The remaining determinant is that of $V$,
of absolute value one. Evaluation of a real cohomology on these
cycles therefore forces all coefficients, including the constant,
to vanish.
\end{proof}

The preceding theorems are statements about real billiard orbits, not
merely formal integer records. They also exclude a circle-valued
coboundary with a continuous transfer function on the regular return
branches: an almost-everywhere identity then extends to each regular
periodic point by continuity and full support of the section measure.
A transfer function assumed only measurable requires a separate
regularity theorem. Similarly, a covariance conclusion requires a
proved variance--coboundary criterion in the relevant function space.
If that criterion and a continuous covariance family are established,
Corollary~\ref{cor:real-period-rank} gives positive definiteness and,
on $I$, a uniform positive lower eigenvalue. Existence and continuity
of that covariance are not consequences of the determinant alone.
Compact periodic separation is not a quantitative operator contraction
or a returned temporal nonintegrability estimate.

\section{Quantitative separation by logarithmic-length periods}
\label{sec:quantitative-periods}
The complete periodic annihilator was determined in
Theorem~\ref{thm:joint-periodic}. We now estimate the period increments
uniformly in both the radius and the number of collisions. This gives an
explicit separation at large roof frequencies, using three actual
periodic words. Throughout this section $g=\sqrt3-2R$, $d=1/2-R$,
and $e$, $\ell$ and $u$ are the functions in
\eqref{eq:excursion-lengths}. The coordinates below are heights, not
arc coordinates.

Define the half action on $[-d,0]^m$ by
\begin{equation}\label{eq:half-action-v4}
 W_{m,R}(y)=e_R(y_1)-g/2+
       \sum_{j=1}^{m-1}\bigl(\ell_R(y_j,y_{j+1})-g\bigr)+u_R(y_m).
\end{equation}
Reflection symmetry and uniqueness of the full minimizer imply
$E_m(R)=2\min W_{m,R}$. Denote its minimizing coordinates by
$y^{(m,R)}_1,\ldots,y^{(m,R)}_m$.

\begin{lemma}[Uniform Hessian and endpoint estimates]
\label{lem:uniform-half-hessian}
On the entire height box, the Hessian $H=\nabla^2 W_{m,R}$ is
tridiagonal and satisfies
\begin{equation}\label{eq:half-hessian-bounds}
 3I_m\le H\le12I_m,\qquad H_{jj}\le10,\qquad
 \frac12\le-H_{j,j+1}\le2\quad(1\le j<m).
\end{equation}
The last condition is empty when $m=1$. For all $1\le j\le m$,
\begin{equation}\label{eq:height-two-sided}
 \frac3{440}\,20^{-(j-1)}
 \le -y^{(m,R)}_j
 \le \frac1{21}\left(\frac7{10}\right)^{j-1}.
\end{equation}
All constants are independent of $m$ and $R\in I$.
\end{lemma}
\begin{proof}
Write $x(y)=\sqrt{R^2-y^2}$, $a=D/2-x(y)$,
$h=D-x(y)-x(z)$, $v=z-y$, with $D=\sqrt3$.
The elementary bounds on the whole box are
\[
 \begin{gathered}
 x>11/25,\quad |x'|\le5/44,\quad79/200<a<43/100,\\
 79/100<h<86/100,\quad e<11/25,\quad\ell<9/10.
 \end{gathered}
\]
Also $u''=R^2/x^3\ge1/R\ge100/47$ and
$u''\le55225/21296<3$. For the norm of a vector-valued function,
the Hessian is its transverse positive semidefinite quadratic form
plus the Hessian of the vector dotted with its unit direction.
Consequently
\[
 e''\ge(a/e)u''>\frac{79}{88}\frac{100}{47}>\frac32,
 \qquad
 \nabla^2\ell\ge(h/\ell)\operatorname{diag}(u''(y),u''(z))
 >\frac32 I_2.
\]
Each variable of $W_m$ occurs in two such terms, with $u''>3/2$
at the terminal endpoint. This proves the lower bound $H\ge3I_m$,
including $m=1$.

The same norm formula gives
\[
 \ell_{yy},\ell_{zz}
 \le\frac{1+(5/44)^2}{79/100}+\frac{55225}{21296}<4,
 \qquad
 e''\le\frac{1+(5/44)^2}{79/200}+\frac{55225}{21296}<6.
\]
An explicit mixed derivative is
\begin{equation}\label{eq:mixed-length-hessian}
 -\ell_{yz}
 =\frac{(h+vy/x(y))(h-vz/x(z))}{\ell^3}.
\end{equation}
Each numerator factor lies between $78/100$ and $87/100$.
Using $79/100<\ell<9/10$ proves $1/2<-\ell_{yz}<2$.
The diagonal bounds for $W_m$ are therefore $10$ at its first
endpoint, $8$ in its interior, and $7$ at its last endpoint; when
$m=1$ the bound is $9$. Absolute row sums are at most $12$.
The symmetric matrix bound $H\le12I_m$ follows, proving
\eqref{eq:half-hessian-bounds}.

At zero the gradient is $a_0 e_1$, where
$a_0=d/e_R(0)$ satisfies $3/44\le a_0\le1/7$.
Integrate the Hessian on the segment from zero to the minimizer and
write the resulting matrix as $A$. Stationarity gives
$Ay^{(m,R)}=-a_0e_1$. The matrix $B=10I_m-A$ is entrywise
nonnegative and tridiagonal. Its eigenvalues lie in $[-2,7]$, so
$\|B/10\|_2\le7/10$. Hence
\[
 A^{-1}=\frac1{10}\sum_{k\ge0}(B/10)^k.
\]
For the $(j,1)$ entry every power before $j-1$ vanishes. The
nearest-neighbor product at power $j-1$ is at least $20^{-(j-1)}$.
Entrywise nonnegativity and the norm bound yield
\[
 \frac1{10}20^{-(j-1)}
 \le (A^{-1})_{j1}
 \le\frac1{10}\sum_{k\ge j-1}(7/10)^k
 =\frac13(7/10)^{j-1}.
\]
Multiplication by the two bounds on $a_0$ proves
\eqref{eq:height-two-sided}. The matrix is an averaged Hessian of
the actual nonlinear action; no quadratic approximation of the
minimizing orbit has been made.
\end{proof}

\begin{theorem}[Two-sided excess increments]
\label{thm:quantitative-excess}
For every $R\in I$ and $m\ge1$,
\begin{equation}\label{eq:quantitative-excess}
 \frac1{10000}\,400^{-m}
 \le E_{m+1}(R)-E_m(R)
 \le\frac1{300}\left(\frac{49}{100}\right)^{m-1}.
\end{equation}
In particular $E_m$ converges uniformly on $I$ to a continuous
function $E_\infty$, with
\begin{equation}\label{eq:uniform-excess-limit}
 0<E_\infty(R)-E_m(R)
 \le\frac1{153}\left(\frac{49}{100}\right)^{m-1}.
\end{equation}
\end{theorem}
\begin{proof}
For a minimizer $y$ of $W_m$, append the coordinate zero. With
$K(y,z)=\ell(y,z)-g$, the resulting change in half action is
\[
 K(y_m,0)-u(y_m)
 =\sqrt{(g+u(y_m))^2+y_m^2}-(g+u(y_m))
 \le\frac{y_m^2}{2g}.
\]
It follows that $E_{m+1}-E_m\le y_m^2/g$. The upper bound in
\eqref{eq:height-two-sided} and $g>79/100$ give the upper
coefficient $100/34839<1/300$ in \eqref{eq:quantitative-excess}.

Conversely, truncate a minimizer $z$ of $W_{m+1}$ after its first
$m$ coordinates. The removed contribution is
\[
 K(z_m,z_{m+1})+u(z_{m+1})-u(z_m)\ge2u(z_{m+1}),
\]
by \eqref{eq:kernel-coercivity}. Thus
\[
 E_{m+1}-E_m\ge4u(z_{m+1})\ge\frac{2z_{m+1}^2}{R}
 \ge\frac9{45496}\,400^{-m}>\frac1{10000}\,400^{-m}.
\]
Here $u(z)=z^2/(R+x(z))\ge z^2/(2R)$ and
\eqref{eq:height-two-sided} were used. Both comparisons concern
minima of the exact actions. Summing the geometric upper bound
proves \eqref{eq:uniform-excess-limit}. Each $E_m$ is continuous
by Lemma~\ref{lem:excursion-family}; the uniform limit is continuous.
\end{proof}

Let $P_0$ denote the long normal orbit and $P_m$ the $m$th excursion.
For a periodic word $P$ write
\[
 Z_P(u,s,\beta,c;R)
 =\exp i\bigl(u\cdot K_P+sN_P+\beta T_P-ch_P\bigr).
\]
Put $a=\log(100/49)$ and, for $b\ge1$, define
\begin{equation}\label{eq:frequency-period-budget}
 m(b)=1+\left\lceil\frac{\log b}{a}\right\rceil,
 \qquad \gamma=\frac{\log400}{a}-1.
\end{equation}

\begin{theorem}[Explicit large-roof phase separation]
\label{thm:log-period-separation}
For $|\beta|=b\ge1$, all $R\in I$, and arbitrary torus
frequencies $u,s,c$, the three true periods $P_0,P_{m(b)},P_{m(b)+1}$
satisfy
\begin{equation}\label{eq:log-period-separation}
 \max_P|Z_P-1|
 \ge\frac{b\,400^{-m(b)}}{15000\pi}
 \ge\frac{b^{-\gamma}}{2400000000\pi}.
\end{equation}
The exponent is $\gamma=\log400/\log(100/49)-1$.
For a whole band $1\le|\beta|\le B$, the fixed three periods
$P_0,P_{m(B)},P_{m(B)+1}$ suffice, with lower bound
$400^{-m(B)}/(15000\pi)$. Their longest physical word has
$2m(B)+3$ collisions. All assertions are uniform in the radius.
\end{theorem}
\begin{proof}
The actual counts and lengths give the exact identity
\begin{equation}\label{eq:three-period-cancellation}
 Z_{P_{m+1}}\overline{Z_{P_m}}\overline{Z_{P_0}}
       =\exp\bigl(i\beta(E_{m+1}-E_m)\bigr).
\end{equation}
The winding, physical count and induced period cancel. With $m=m(b)$,
\eqref{eq:quantitative-excess} implies
\[
 \frac{b\,400^{-m}}{10000}
 \le b(E_{m+1}-E_m)\le\frac1{300}.
\]
For $|v|\le\pi$, $|e^{iv}-1|\ge2|v|/\pi$. The distance
from one of a product of three unit complex numbers is at most
the sum of their individual distances from one. Applying these
facts to \eqref{eq:three-period-cancellation} proves the first
bound in \eqref{eq:log-period-separation}. Since
$m(b)\le2+\log b/a$, the second bound follows. For the band
assertion use $m(B)$; the upper phase bound still holds, while
$b\ge1$ supplies the printed lower bound. The number of collisions
was proved in Lemma~\ref{lem:excursion-family}.
\end{proof}

\begin{corollary}[An approximate-phase lower bound]
\label{cor:approximate-phase-bound}
Let $q:Y_R^*\to\mathbb C$ have modulus one and be continuous in
neighborhoods of all states in the three cycles used above. Suppose
$q$ is measurable elsewhere and let
\[
 D(q)=\mathop{\rm ess\,sup}_{x\in Y_R^*}
 \left|e^{i(u\cdot\kappa_R^*(x)+s r_R^*(x)+\beta\tau_R^*(x))}
                 q(F_R^*x)-e^{ic}q(x)\right|.
\]
For $b=|\beta|\ge1$ and $m=m(b)$,
\begin{equation}\label{eq:approximate-phase-bound}
 D(q)\ge\frac{b\,400^{-m}}{10000\pi(m+1)}.
\end{equation}
\end{corollary}
\begin{proof}
The return branches through these cycles are regular, and their
section measure has full support. Continuity therefore extends the
essential bound to every state of each selected cycle. Multiplication
around a cycle of induced period $h$ telescopes the unit phases and
gives $|Z_P-1|\le hD(q)$. In
\eqref{eq:three-period-cancellation} the sum of the three periods is
$1+m+(m+1)=2(m+1)$. The lower sine bound in the preceding proof
therefore gives
$2b(E_{m+1}-E_m)/\pi\le2(m+1)D(q)$. Apply the lower bound
in \eqref{eq:quantitative-excess}.
\end{proof}

The estimates are quantitative statements about the specified mechanical
return map, and in particular about continuous circle-valued approximate
phases. They do not assert regularity of a merely measurable transfer
function, a lower bound for a spectral eigenfunction on these orbits,
or a transfer-operator resolvent estimate. Those are distinct analytical
steps when using \eqref{eq:approximate-phase-bound} in a high-frequency
argument. The raw characteristic-function tail is not replaced by the
period discrepancy.

\section{Uniform phase separation and positive-measure coercivity}
\label{sec:periodic-coercivity}
We strengthen the frequency estimate by comparing consecutive exact
minimizers. The choice of a period will depend on the frequency and
possibly on the radius. This dependence is essential: a fixed triple
of periods is not asserted to separate the entire unbounded frequency
axis by a positive constant.

Retain $K_R(y,z)=\ell_R(y,z)-g$ and $W_{m,R}$ from
\eqref{eq:half-action-v4}. Write
\[
 a_m(R)=-y_m^{(m,R)}>0,\qquad
 \Delta_m(R)=E_{m+1}(R)-E_m(R)>0.
\]
The symbol $a_m$ denotes a terminal height; it is unrelated to the
logarithmic constant in \eqref{eq:frequency-period-budget}.

\subsection{Comparison of successive exact minima}
\begin{lemma}[The terminal value function]
\label{lem:terminal-response}
For $y\in[-d,0]$ put
\[
 V_R(y)=\min_{z\in[-d,0]}\{K_R(y,z)+u_R(z)\},\qquad
 \delta_R(y)=V_R(y)-u_R(y).
\]
The minimizer $f_R(y)$ is unique. It belongs to $(-d,0)$ for $y<0$
and equals zero for $y=0$. The functions are twice continuously
differentiable, with one-sided derivatives at the endpoints, and
\begin{equation}\label{eq:terminal-response}
 f_R'(y)\ge\frac1{14},\qquad
 |f_R(y)|\ge\frac{|y|}{14},\qquad
 |\delta_R''(y)|\le2,\qquad |\delta_R'(y)|\le2|y|.
\end{equation}
These estimates hold on the entire height interval, uniformly in $R$.
\end{lemma}
\begin{proof}
The second derivative in $z$ is $K_{zz}+u''>0$. At $z=-d$,
\[
 K_z(y,-d)=\frac{h(y,-d)(-d)/x_R(-d)-d-y}{\ell_R(y,-d)}<0,
 \qquad u_R'(-d)<0.
\]
At $z=0$ and $y<0$ the derivative is $-y/\ell_R(y,0)>0$.
Consequently the minimizer is interior and unique. For $y=0$,
\eqref{eq:kernel-coercivity} gives the unique minimizer zero.
The implicit function theorem applies to the stationarity equation,
including its smooth extension near the endpoints, and gives
\[
 f_R'(y)=\frac{-K_{yz}(y,f_R(y))}
                    {K_{zz}(y,f_R(y))+u_R''(f_R(y))}
 \ge\frac{1/2}{4+3}=\frac1{14}.
\]
Here the individual segment bounds established in the proof of
Lemma~\ref{lem:uniform-half-hessian} were used. Integration from $y$
to zero proves the second inequality of \eqref{eq:terminal-response}.

The envelope formula is
\[
 V_R''(y)=K_{yy}(y,f_R(y))-
       \frac{K_{yz}(y,f_R(y))^2}{K_{zz}(y,f_R(y))+u_R''(f_R(y))}.
\]
The Hessian of $K_R(y,z)+u_R(z)$ dominates $(3/2)I_2$.
Its Schur complement therefore is at least $3/2$: evaluate its
quadratic form on $(1,t)$ and minimize over $t$. Also
$V_R''\le K_{yy}\le4$. Since $100/47\le u_R''<3$, subtraction
shows $|\delta_R''|\le2$. The first derivatives of $K_R$ and $u_R$
vanish at the origin, so $\delta_R'(0)=0$. Integration proves the
last inequality. All estimates concern the nonlinear action itself.
\end{proof}

\begin{proposition}[Consecutive terminal heights]
\label{prop:terminal-comparison}
For every $m\ge1$ and $R\in I$,
\begin{equation}\label{eq:terminal-comparison}
 a_{m+1}(R)\ge\frac3{70}a_m(R).
\end{equation}
\end{proposition}
\begin{proof}
Let $y$ minimize $W_{m,R}$ and let $(X,z)$ minimize $W_{m+1,R}$,
where $X$ consists of the first $m$ coordinates. Eliminating the
last coordinate gives exactly
\[
 \min_z W_{m+1,R}(X,z)=W_{m,R}(X)+\delta_R(X_m),
 \qquad z=f_R(X_m).
\]
All minimizing coordinates are interior by
Lemma~\ref{lem:excursion-family}. Stationarity and the fundamental
theorem of calculus thus yield
\[
 A(X-y)=-\delta_R'(X_m)e_m,\qquad
 A=\int_0^1\nabla^2W_{m,R}(y+t(X-y))\dd t\ge3I_m.
\]
The segment stays in the height box. In particular
$0<(A^{-1})_{mm}\le1/3$. Taking the last coordinate and applying
Lemma~\ref{lem:terminal-response}, we obtain
\[
 |y_m|\le |X_m|+\frac13|\delta_R'(X_m)|
             \le\frac53|X_m|.
\]
Finally $|z|=|f_R(X_m)|\ge|X_m|/14$. These two inequalities give
\eqref{eq:terminal-comparison}. No monotonicity of the terminal
heights, and no quadratic replacement of $W_{m,R}$, is assumed.
\end{proof}

\begin{corollary}[A lower ratio for the excess increments]
\label{cor:increment-ratio}
For every $m\ge1$ and $R\in I$,
\begin{equation}\label{eq:increment-ratio}
 \Delta_{m+1}(R)\ge10^{-5}\Delta_m(R).
\end{equation}
\end{corollary}
\begin{proof}
The append and truncate comparisons in the proof of
Theorem~\ref{thm:quantitative-excess} give
\[
 \Delta_m\le\frac{a_m^2}{g},\qquad
 \Delta_{m+1}\ge\frac2R a_{m+2}^2.
\]
Applying \eqref{eq:terminal-comparison} twice, and using
$g>79/100$ and $R\le47/100$, yields
\[
 \Delta_{m+1}\ge\frac{2g}{R}\left(\frac3{70}\right)^4\Delta_m
 >\frac{158}{47}\left(\frac3{70}\right)^4\Delta_m
 >\frac{\Delta_m}{100000}.
\]
This is a comparison between adjacent increments, stronger for the
present purpose than separate exponential upper and lower bounds.
\end{proof}

\subsection{A frequency-dependent triple and a fixed finite menu}
Put
\begin{equation}\label{eq:uniform-phase-constants}
 \eta_0=\frac1{4000000},\qquad
 \sigma_0=\frac1{6000000\pi},\qquad
 M(b)=m(b)+1,
\end{equation}
where $m(b)$ is defined in \eqref{eq:frequency-period-budget}.

\begin{theorem}[Uniform positive phase separation]
\label{thm:uniform-phase-separation}
For every $R\in I$ and $b\ge1$ the integer
\begin{equation}\label{eq:adaptive-period}
 j(R,b)=\min\{j\ge1:b\Delta_j(R)\le1/2\}
\end{equation}
is well defined and satisfies $j(R,b)\le m(b)$. If $|\beta|=b$,
then, for arbitrary $u,s,c$,
\begin{equation}\label{eq:uniform-phase-separation}
 \max_{P\in\{P_0,P_{j(R,b)},P_{j(R,b)+1}\}}
       |Z_P(u,s,\beta,c;R)-1|\ge\sigma_0.
\end{equation}
The longest of these physical words has at most $2m(b)+3$ collisions.
For the entire band $1\le|\beta|\le B$, the fixed menu
\begin{equation}\label{eq:fixed-period-menu}
 \mathcal P_B(R)=\{P_0(R),P_1(R),\ldots,P_{m(B)+1}(R)\}
\end{equation}
satisfies the same lower bound for its maximum, uniformly in $R$.
It contains $m(B)+2$ periods.
\end{theorem}
\begin{proof}
The upper bound of \eqref{eq:quantitative-excess} gives
$b\Delta_{m(b)}\le1/300$. Thus the defining set in
\eqref{eq:adaptive-period} is nonempty. If $j(R,b)=1$, its lower
bound gives $b\Delta_1\ge\Delta_1\ge1/4000000=\eta_0$.
If $j(R,b)>1$, minimality and \eqref{eq:increment-ratio} give
\[
 b\Delta_{j(R,b)}>10^{-5}/2\ge\eta_0.
\]
In either case
\begin{equation}\label{eq:selected-phase-window}
 \eta_0\le b\Delta_{j(R,b)}\le1/2.
\end{equation}
The exact three-period identity \eqref{eq:three-period-cancellation}
therefore has distance at least $2\eta_0/\pi$ from one. The distance
of a product of three unit complex numbers from one is at most the
sum of their three distances. This proves
\eqref{eq:uniform-phase-separation}, since
$2\eta_0/(3\pi)=\sigma_0$. The physical and induced counts are
those of Lemma~\ref{lem:excursion-family}; hence the stated word
length. Finally $m(b)\le m(B)$ on the band, so the selected triple
is contained in \eqref{eq:fixed-period-menu}.

The index selection uses a first crossing, not monotonicity of
$\Delta_j$. It need not be continuous in $R$ or $b$. The menu has
fixed indices and consists of analytically varying true periods,
which is the asserted parameter-uniform conclusion.
\end{proof}

\begin{corollary}[Logarithmic uniform-norm coercivity]
\label{cor:logarithmic-coercivity}
Let $q$ be circle-valued and measurable on $Y_R^*$, and continuous
in neighborhoods of every state of the triple in
\eqref{eq:uniform-phase-separation}. With $D(q)$ as in
Corollary~\ref{cor:approximate-phase-bound}, for $b=|\beta|\ge1$,
\begin{equation}\label{eq:logarithmic-coercivity}
 D(q)\ge\frac1{4000000\pi(j(R,b)+1)}
       \ge\frac1{4000000\pi M(b)}.
\end{equation}
\end{corollary}
\begin{proof}
Continuity and full support extend the essential bound to the
selected regular periodic states. Telescoping around each period
bounds its phase discrepancy by its induced length times $D(q)$.
The three induced lengths sum to $2(j(R,b)+1)$. Apply the product
identity and \eqref{eq:selected-phase-window} to obtain
\[
 2\eta_0/\pi\le2(j(R,b)+1)D(q).
\]
This proves \eqref{eq:logarithmic-coercivity}. The modulus-one
condition is used in telescoping and is not a consequence of an
arbitrary operator normalization.
\end{proof}

\subsection{Regular tubes of the actual return map}
We use the Euclidean metric in local $(\alpha,p)$ coordinates. The
following uniformity concerns the specified periodic family, not all
branches of the induced map.

\begin{lemma}[Uniform return neighborhoods of the period family]
\label{lem:uniform-period-tubes}
There are constants $0<\rho_0<1$, $L\ge2$ and $B_0\ge1$, independent
of $R$ and $m$, with the following properties. For every selected
state $x_P$ of $P_0$ or any $P_m$, the ball $B(x_P,\rho_0)$ lies
in $Y_R^*$ and in one regular first-return branch. On this ball
$r_R^*$ and $\kappa_R^*$ are constant, $r_R^*\in\{2,3\}$, and
\[
 \operatorname{Lip}(F_R^*)\le L,\qquad
 \operatorname{Lip}(\tau_R^*)\le B_0.
\]
If $P$ has induced period $h$ and $0<\rho\le\rho_0L^{-h}$, every
$x\in B(x_P,\rho)$ follows the same $h$ return branches as $x_P$,
and, writing $F=F_R^*$,
\begin{align}\label{eq:period-tube-bounds}
 d(F^k x,F^k x_P)&\le L^k\rho &&(0\le k\le h),\\
 |T_{h,R}(x)-T_{h,R}(x_P)|&\le2B_0L^h\rho,\\
 \nu_R^*(B(x_P,\rho))&=\rho^2/(4c_*).
\end{align}
\end{lemma}
\begin{proof}
We verify that the family stays a uniform distance from every local
singularity relevant to a single return. The estimates in
Lemma~\ref{lem:excursion-family} hold for all its contact heights:
$-2d/5<y<0$, the long-flight third-disk clearances are at least
$9/500$, and the incidence cosine at $A$ exceeds $1/25$.
At the facing contacts on $O,C$ it is bounded below by a positive
constant from the horizontal-component inequalities in that proof.
End flights have a uniform positive distance from every third disk:
the opposite facing center is horizontally separated by $D/2$,
$B$ is vertically separated, and the remaining lattice centers obey
the enumerated uniform separation bounds. All flight lengths are
bounded above and below by positive constants.

Return membership has uniform margins as well. The selected contacts
on $O$ satisfy
\[
 |\alpha-\pi/6|\le1/22<3/50,\qquad
 |p|\le1/22+2/79<3/20.
\]
The unselected contacts on $C$ and $A$ have angles near $7\pi/6$
and equal to $2\pi/3$, respectively, uniformly separated from $U$
and the four squares in $Y_R$. The long normal orbit has the same
strict margins. Each selected-to-selected path consequently contains
two physical flights, or three through $C,A,O$.

For clarity, compactness is applied only to these bounded-length
paths. The parameters $R$ and their contact coordinates range in a
compact box. Every limiting path satisfies the non-grazing,
nonocclusion and return-membership margins just established. In a
neighborhood of such a path the selected entry roots are simple;
the physical collision map and flight time are analytic in the
initial state and parameter. There are only the finitely many
center itineraries just listed, up to lattice translation. A finite
cover of their compact closures therefore gives a common neighborhood
radius and uniform first-derivative bounds for the actual two- or
three-flight first returns. Decrease the radius so its coordinate
balls lie in $Y_R^*$ and contain no change of lattice labels; enlarge
the derivative bounds to $L\ge2$, $B_0\ge1$. This proves the first
part without taking a derivative of an arbitrarily long return.

Starting in a ball of radius at most $\rho_0L^{-h}$, induction gives
the first inequality in \eqref{eq:period-tube-bounds} and preserves
the same return decisions at each step. Summing the roof variations
gives $B_0\rho\sum_{k<h}L^k\le2B_0L^h\rho$. Finally
$\dd\nu_R^*=(4\pi c_*)^{-1}\dd\alpha\dd p$ on the ball, whose
Euclidean area is $\pi\rho^2$. This proves the last equality.
\end{proof}

\begin{theorem}[Positive-measure coercivity for regular phases]
\label{thm:positive-measure-coercivity}
Fix $0<\alpha\le1$, $H\ge1$ and $1\le p<\infty$. Suppose $q$ is
measurable and circle-valued on $Y_R^*$ and has an $\alpha$-H\"older
representative, of constant at most $H$, on the balls of radius
$\rho_0$ around all states of the selected triple. Define
\[
 \mathcal D_q(x)=e^{i(u\cdot\kappa_R^*(x)+s r_R^*(x)+\beta\tau_R^*(x))}
                      q(F_R^*x)-e^{ic}q(x).
\]
For $b=|\beta|\ge1$, put
\begin{equation}\label{eq:coercivity-radius}
 \rho(b,H)=L^{-M(b)}
 \min\left\{\frac{\rho_0}{2},\frac{\sigma_0}{8bB_0},
                       \left(\frac{\sigma_0}{8H}\right)^{1/\alpha}\right\}.
\end{equation}
Then the actual invariant return probability satisfies
\begin{equation}\label{eq:positive-measure-coercivity}
 \|\mathcal D_q\|_{L^p(\nu_R^*)}
 \ge\frac{\sigma_0}{2M(b)}
             \left(\frac{\rho(b,H)^2}{4c_*}\right)^{1/p}.
\end{equation}
All geometric constants are uniform over $R\in I$. No independence
of the returns is assumed.
\end{theorem}
\begin{proof}
Select from the triple a period $P$ with $|Z_P-1|\ge\sigma_0$.
Let its selected state be $x_P$ and its induced period be $h\le M(b)$.
Write $\varphi=u\cdot\kappa_R^*+s r_R^*+\beta\tau_R^*$, $F=F_R^*$
and $S_h\varphi=\sum_{k<h}\varphi\circ F^k$. The lattice and count
terms are constant along the corresponding branches in the tube.
Thus for $x\in B(x_P,\rho(b,H))$, Lemma~\ref{lem:uniform-period-tubes}
gives
\[
 |S_h\varphi(x)-S_h\varphi(x_P)|\le2bB_0L^h\rho(b,H).
\]
At $x_P$, the modulus of
$e^{iS_h\varphi(x_P)}q(F^h x_P)-e^{ich}q(x_P)$ is $|Z_P-1|$.
Since $F^h x_P=x_P$, the H\"older bounds at the two endpoints and
\eqref{eq:coercivity-radius} imply, throughout the smaller ball,
\begin{align*}
 \left|e^{iS_h\varphi(x)}q(F^h x)-e^{ich}q(x)\right|
 &\ge\sigma_0-2bB_0L^h\rho(b,H)-2H(L^h\rho(b,H))^\alpha\\
 &\ge\sigma_0/2.
\end{align*}
Choose the indicated representative on these balls; changes on null
sets do not affect the following integrals. Telescoping $h$ times
bounds the left side by $\sum_{k<h}|\mathcal D_q(F^k x)|$ almost
everywhere. Minkowski's inequality on the smaller ball and invariance
of $\nu_R^*$ give
\[
 \frac{\sigma_0}{2}\nu_R^*(B(x_P,\rho(b,H)))^{1/p}
 \le\sum_{k<h}\|\mathcal D_q\circ F^k\|_{L^p(B(x_P,\rho(b,H)),\nu_R^*)}
 \le h\|\mathcal D_q\|_{L^p(\nu_R^*)}.
\]
Use the exact ball mass and $h\le M(b)$ to conclude. The estimate
uses measure preservation, not disjointness or probabilistic
independence of the iterates of the ball.
\end{proof}

\begin{corollary}[A polynomial budget for the phase regularity]
\label{cor:polynomial-regularity-budget}
Fix $\alpha,p,H_0$ as above and $\zeta\ge0$. If the H\"older constants
satisfy $H\le H_0b^\zeta$, set
\begin{equation}\label{eq:coercivity-exponent}
 \kappa=\frac{\log L}{\log(100/49)}+
                       \max\{1,\zeta/\alpha\}.
\end{equation}
There is $C>0$, depending only on these fixed constants and the
mechanical family, such that
\begin{equation}\label{eq:polynomial-coercivity}
 \|\mathcal D_q\|_{L^p(\nu_R^*)}
             \ge C\frac{b^{-2\kappa/p}}{1+\log b}\qquad(b\ge1).
\end{equation}
\end{corollary}
\begin{proof}
Writing $a=\log(100/49)$, we have $M(b)\le3+\log b/a$ and hence
$L^{-M(b)}\ge L^{-3}b^{-\log L/a}$. The minimum in
\eqref{eq:coercivity-radius} is at least
\[
 \min\left\{\frac{\rho_0}{2},\frac{\sigma_0}{8B_0},
           \left(\frac{\sigma_0}{8H_0}\right)^{1/\alpha}\right\}
                     b^{-\max\{1,\zeta/\alpha\}}.
\]
Substitution into \eqref{eq:positive-measure-coercivity} proves the
claim with a positive constant independent of $R$ and $b$.
\end{proof}

\subsection{The operator interface}
Theorems~\ref{thm:uniform-phase-separation} and
\ref{thm:positive-measure-coercivity} have different hypotheses.
The former is an unconditional estimate for the specified mechanical
periods. The latter is an estimate for a stated regularity class of
circle-valued phases, in the actual invariant measure. In particular
it supplies a positive-measure obstruction rather than a conclusion
obtained by evaluating an unspecified measurable representative at
isolated periodic points.

To apply \eqref{eq:polynomial-coercivity} to a transfer operator,
one must construct from its approximate spectral vectors phases in
that class, bound the H\"older constant, and compare their physical
$L^p$ defect with the operator defect. Such a comparison would
contradict an error smaller than the right side of
\eqref{eq:polynomial-coercivity}. It is not proved merely by writing
an eigenvalue equation. In particular a division by the modulus of a
spectral vector requires a separate lower-modulus estimate.
General Liv\v{s}ic regularity results for Markov systems
\cite{BHN} concern precise dynamical and cohomological hypotheses;
they do not supply a uniform quantitative regularity bound for these
approximate phases without checking those hypotheses and the defect
comparison. No such comparison is assumed here for the full induced
operator. The common operator realization and the all-branch raw
residual estimates in the final section remain the requirements for
the original joint density theorem.

\section{One-step phase coercivity and the operator criterion}
\label{sec:operator-bridge}
The positive-measure estimate in Section~\ref{sec:periodic-coercivity}
uses an entire periodic tube. There is a stronger estimate if one first
locates a large one-step defect on the period, and only then thickens
that single state. It avoids the exponentially shrinking tube of the
whole period. We retain the constants $\rho_0,L,B_0$ of
Lemma~\ref{lem:uniform-period-tubes}, and set $r_0=\rho_0/(2L)$.

\begin{theorem}[One-step thickening of the periodic defect]
\label{thm:one-step-coercivity}
Fix $0<\alpha\le1$ and $1\le p<\infty$. Let $q$ be circle-valued and
measurable, with an $\alpha$-H\"older representative of constant $H\ge1$
on the $\rho_0$ balls around the states of the triple selected in
Theorem~\ref{thm:uniform-phase-separation}. For $b=|\beta|\ge1$ put
\[
 d_b=\frac1{4000000\pi M(b)},\qquad
 H_b=H(1+L^\alpha),\qquad
 r_b=\min\{r_0,d_b/(4bB_0),(d_b/(4H_b))^{1/\alpha}\}.
\]
For the actual one-step defect $\mathcal D_q$ defined in
Theorem~\ref{thm:positive-measure-coercivity},
\begin{equation}\label{eq:one-step-coercivity}
 \|\mathcal D_q\|_{L^p(\nu_R^*)}
       \ge\frac{d_b}{2}\left(\frac{r_b^2}{4c_*}\right)^{1/p}.
\end{equation}
If $H\le H_0b^\zeta$ with fixed $H_0\ge1$ and $\zeta\ge0$, this gives
\begin{equation}\label{eq:sharper-phase-power}
 \|\mathcal D_q\|_{L^p(\nu_R^*)}
 \ge C\,b^{-\lambda}(1+\log b)^{-\eta},\qquad
 \lambda=\frac2p\max\{1,\zeta/\alpha\},\quad
 \eta=1+\frac2{\alpha p}.
\end{equation}
The constants are uniform in the radius and all torus frequencies.
\end{theorem}
\begin{proof}
Telescoping the three exact periods as in
Corollary~\ref{cor:logarithmic-coercivity} shows that at some state $x_0$
of the selected triple the one-step defect has modulus at least $d_b$.
This is initially a statement about the specified continuous
representatives at finitely many states: the sum of the three induced
periods is at most $2M(b)$, while their product discrepancy is at least
$2/(4000000\pi)$.

On $B(x_0,r_0)$ the return has one fixed lattice label and one fixed
physical count, and $F_R^*$ maps that ball into the $\rho_0$ ball about
the successor of $x_0$. For $x$ in that ball,
\[
 |\mathcal D_q(x)-\mathcal D_q(x_0)|
 \le H_b|x-x_0|^\alpha+bB_0|x-x_0|.
\]
Each term on the right is at most $d_b/4$ on the ball of radius
$r_b$. On that ball the defect
is consequently at least $d_b/2$. Its actual invariant measure is
$r_b^2/(4c_*)$. Integration proves \eqref{eq:one-step-coercivity}.
Changes of representatives on null sets do not alter the integral.

Finally $M(b)\le3+\log b/\log(100/49)$ and
$H_b\le C b^\zeta$. The minimum defining $r_b$ is bounded
below by a positive constant times
$b^{-\max\{1,\zeta/\alpha\}}(1+\log b)^{-1/\alpha}$.
Substitution proves \eqref{eq:sharper-phase-power}. Unlike the
whole-period tube estimate, its exponent contains no $\log L$ term.
The regularity premise on $q$ remains essential.
\end{proof}

The next theorem specifies exactly what is required to turn this
physical coercivity into an operator resolvent. The hypothesis is a
reconstruction of a phase from an approximate spectral vector; it is
not silently identified with an eigenvalue equation.

\begin{theorem}[A quantitative phase-reconstruction criterion]
\label{thm:phase-resolvent-criterion}
For each $R,u,s,\beta$, let $\mathcal L_{R,u,s,\beta}$ be a bounded
operator on a Banach space $\mathcal B_R$. Fix the constants
$\alpha,p,H_0,\zeta$ of Theorem~\ref{thm:one-step-coercivity}, and let
$A\ge1$, $v\ge0$, $\theta>0$. Suppose that for every $b=|\beta|\ge1$,
$|z|=1$, and $h\in\mathcal B_R$ with $\|h\|=1$ and
$e=\|(z-\mathcal L_{R,u,s,\beta})h\|\le1$, one can construct a
circle-valued $q$ with the local H\"older bound $H_0b^\zeta$ and
\begin{equation}\label{eq:phase-reconstruction}
 \|\mathcal D_q\|_{L^p(\nu_R^*)}\le A b^v e^\theta,
 \qquad e^{ic}=z.
\end{equation}
Suppose also that $z-\mathcal L_{R,u,s,\beta}$ is Fredholm of index
zero. Then it is invertible and
\begin{equation}\label{eq:phase-resolvent}
 \|(z-\mathcal L_{R,u,s,\beta})^{-1}\|
 \le C b^{(\lambda+v)/\theta}
                (1+\log b)^{\eta/\theta},
\end{equation}
with $\lambda,\eta$ as in \eqref{eq:sharper-phase-power}.
\end{theorem}
\begin{proof}
Combining \eqref{eq:sharper-phase-power} and
\eqref{eq:phase-reconstruction} yields
\[
 e\ge (C_0/A)^{1/\theta}
       b^{-(\lambda+v)/\theta}(1+\log b)^{-\eta/\theta}
\]
when $e\le1$. Decreasing the positive constant to at most one makes
the same lower bound valid when $e>1$. Homogeneity gives a uniform
lower bound for $z-\mathcal L$ on every vector. It is therefore
injective and has closed range. Fredholm index zero makes its cokernel
zero as well, so it is surjective. The lower bound gives
\eqref{eq:phase-resolvent}. The Fredholm assumption cannot be omitted:
an injective operator bounded below need not be onto.
\end{proof}

For the unbounded induced twists, the phase reconstruction
\eqref{eq:phase-reconstruction}, including nonvanishing modulus and
quantitative regularity, is an unverified premise, not a result of
Lemma~\ref{lem:collision-spectral-input}. The latter treats a smooth
collision-space weight near zero. It provides the return-tail theorem
but is not a construction of $\mathcal B_R$ or of the raw high-frequency
operator. The criterion separates the necessary analytical construction
from the now explicit phase lower bound.

\paragraph{Scope of the retained edge calculation.}
The next section uses the original section $Y_R$ in \eqref{eq:Y}, not
$Y_R^*$. Its pointwise physical critical-word calculations are retained,
but an induced count or normalization is not transferred between the
sections without a new return computation.
\section{The exact raw density at a minimal flight length}
Let $L_N(x)=S_N\tau_R(x)$ on the collision section, and put
\[
 g=1-2R,\qquad \zeta_R=\operatorname{arcosh}(1+g/R)>0.
\]
Near the orbit alternating normally between centers $0$ and $a$, choose the initial arc coordinate $y$ in the upward tangent direction and the outgoing tangential velocity $p$. Both vanish on the orbit. Use the same upward signed arc coordinate $y_j$ at its successive endpoints.

\begin{lemma}[Reduced action Hessian]\label{lem:hessian}
The second variation of the length with all contact coordinates free is
\begin{equation}\label{eq:quadratic}
 \frac12\sum_{j=0}^{N-1}
 \left\{\frac{y_j^2+y_{j+1}^2}{R}
             +\frac{(y_{j+1}-y_j)^2}{g}\right\}.
\end{equation}
After solving the internal reflection equations, the endpoint Hessian is
\begin{equation}\label{eq:schur}
 S_N=\begin{pmatrix}A_N&-B_N\\-B_N&A_N\end{pmatrix},\quad
 A_N=\frac{\sinh\zeta_R}{g}\coth(N\zeta_R),\quad
 B_N=\frac{\sinh\zeta_R}{g\sinh(N\zeta_R)}.
\end{equation}
In the physical initial coordinates $(y,p)$, $L_N$ has a nondegenerate minimum $Ng$ and Hessian $H_N$ satisfying
\begin{equation}\label{eq:detH}
             \sqrt{\det H_N}=\sinh(N\zeta_R).
\end{equation}
\end{lemma}
\begin{proof}
At successive facing circles the longitudinal endpoint corrections are $y_j^2/(2R)$ and $y_{j+1}^2/(2R)$. Expanding their distance gives \eqref{eq:quadratic}. Its full Hessian is positive definite, since it is a sum of positive square terms including endpoint squares. The interior stationarity equations are
\[
 y_{j+1}-2(1+g/R)y_j+y_{j-1}=0.
\]
Their solution with given endpoints is
\[
 y_j=\frac{\sinh((N-j)\zeta_R)}{\sinh(N\zeta_R)}y_0
       +\frac{\sinh(j\zeta_R)}{\sinh(N\zeta_R)}y_N.
\]
Substitution into the endpoint gradients yields \eqref{eq:schur}. The nonlinear interior equations are uniquely solvable near zero by their positive definite Hessian. Their solution is the physical reflected path by the first variation and the strict incidence at the normal orbit.

The endpoint first variation gives $p=-\partial_{y_0}L_N$ at the linearized level, so $p=-A_Ny_0+B_Ny_N$. The linear map from $(y_0,p)$ to $(y_0,y_N)$ has determinant $1/B_N$. Thus
\[
 \det H_N=\frac{A_N^2-B_N^2}{B_N^2}
                         =\sinh^2(N\zeta_R).
\]
Positivity is preserved under this nonsingular coordinate change, proving the minimum assertion.
\end{proof}

\begin{theorem}[A physical edge coefficient]\label{thm:edge}
Let $\chi$ be a smooth nonnegative section cutoff supported in the regular alternating itinerary, equal to one near its normal initial point. The pushforward of $\chi\nu$ by $L_N$ has, near $Ng$, the form
\begin{equation}\label{eq:edge}
 f_{N,R}(Ng+s)=\one_{\{s\ge0\}}a_{N,R}(s),\qquad
 a_{N,R}(0)=J_{N,R}:=\frac1{2R\sinh(N\zeta_R)}.
\end{equation}
The function $a_{N,R}$ is smooth on a right neighborhood of zero. For each fixed $N$, the construction and smooth bounds may be made uniform on $I$. No assertion of uniform derivative bounds as $N\to\infty$ is made.
\end{theorem}
\begin{proof}
The section density in initial arc and tangential velocity coordinates is $(4\pi R)^{-1}\dd y\dd p$. The Morse lemma applied to Lemma~\ref{lem:hessian} gives $L_N=Ng+|z|^2/2$. The transformed density $w(z)$ is smooth with
\[
 w(0)=\frac1{4\pi R\sqrt{\det H_N}}.
\]
In polar coordinates its pushforward density is
$\one_{\{s\ge0\}}\int_0^{2\pi}w(\sqrt{2s}\,n_\phi)\dd\phi$ near zero. The angular integral of each odd Taylor term vanishes; Taylor's theorem at arbitrary order gives smoothness in $s\ge0$. Its right value is $2\pi w(0)$, which is \eqref{eq:edge}. Compactness in $R$ and the strictly positive determinant give the fixed-$N$ uniform statement.
\end{proof}

\begin{proposition}[Regular critical words and a uniform edge bound]
\label{prop:general-edge}
Fix a regular physical word with $N$ flights between circular obstacles. Every critical point of its total flight length in the two initial section coordinates is normal at both endpoints. There is at most one such point for the fixed sequence of lattice centers. It is a nondegenerate minimum. If $S=\left(\begin{smallmatrix}A&C\\C&D\end{smallmatrix}\right)$ is its reduced endpoint action Hessian, its section-density jump is
\begin{equation}\label{eq:general-edge}
 J=\frac{|C|}{2R\sqrt{AD-C^2}}.
\end{equation}
There is a fixed explicitly computable constant $C_*$ such that, for every regular word and $N\ge2$,
\begin{equation}\label{eq:edge-decay}
                  0<J\le C_*(47/53)^{N-2}.
\end{equation}
No lower incidence margin for the interior collisions is needed for this coefficient bound. This is an individual-word estimate, not a bound for the sum over words.
\end{proposition}
\begin{proof}
Let $q_j(y_j)$ be arc coordinates and $v_j$ the unit velocity of segment $j$. Its second length variation is
\[
 \frac{c_j}{R}y_j^2+\frac{c_{j+1}}{R}y_{j+1}^2+
 \frac{\{v_j^\perp\cdot(t_{j+1}y_{j+1}-t_jy_j)\}^2}{\ell_j},
\]
where $c_j,c_{j+1}>0$ are its outgoing and incoming incidence cosines. At an interior reflected contact the two cosines agree. Summing gives a positive definite tridiagonal matrix with nonzero off-diagonal entries. Its interior Schur complement $S$ is positive definite and $C\ne0$: eliminating the successive interior variables expresses $C$ as the nonzero product of neighboring entries divided by positive pivots.

The endpoint first variation of the reduced action is $-p_0\dd y_0+p_N\dd y_N$. Moreover $\partial_{p_0}y_N=-1/C\ne0$. Thus the derivative of the physical length in $p_0$ vanishes exactly when $p_N=0$, and its other derivative then vanishes exactly when $p_0=0$. At such a critical point the coordinate change gives $\det H=(AD-C^2)/C^2>0$. The same Morse calculation as in Theorem~\ref{thm:edge} proves \eqref{eq:general-edge}.

There cannot be two such regular points for one word. Regard the polygonal length as a convex function of all contact positions in the product of the closed disks. At a normal endpoint its gradient is $-n$, and at an interior reflected contact it is $-2c_j n_j$. These are strictly inward normals. For a different point $x_j$ of the same strictly convex disk, $n_j\cdot(x_j-q_j)<0$. The convex gradient inequality therefore makes every different tuple have strictly larger action. Hence the normal-to-normal critical tuple is the unique global minimizer of that convex problem. Occlusion by a disk not in the word is excluded separately by physical admissibility, not by this minimization argument.

To prove the uniform bound, set $z_j=c_jy_j$ and choose tangent orientations successively so that the off-diagonal entries become negative. At the normal endpoints $c_0=c_N=1$. Write $t_j=1/\ell_j\le t_*=50/3$. The interior matrix has off-diagonals $-t_j$ and diagonals
\[
         d_j=t_{j-1}+t_j+2/(Rc_j)\ge t_{j-1}+t_j+v_*,
             \qquad v_*=200/47.
\]
Its factorization $D_0(I-K)$ has $K\ge0$ with row sums at most
$q=2t_*/(2t_*+v_*)=47/53$. In the Neumann series, an entry connecting interior indices $1$ and $N-1$ vanishes before power $N-2$. Therefore
\[
 |C|\le t_*^2\frac{q^{N-2}}{v_*(1-q)}.
\]
The reduced quadratic form dominates the two endpoint curvature terms, so $S\ge (100/47)I$. Inserting these estimates in \eqref{eq:general-edge}, and using $R\ge9/20$, gives \eqref{eq:edge-decay} with
\[
 C_* = \frac{47}{90}\frac{t_*^2}{v_*(1-q)}.
\]
Arbitrarily small positive interior incidence only increases the diagonal potential, so it does not weaken the estimate.
\end{proof}

The proposition classifies regular critical contributions and bounds each jump uniformly. Singular itinerary boundaries and the number and arrangement of critical values remain to be controlled. Multiplying \eqref{eq:edge-decay} by an uncontrolled word count would not prove a summable global correction.

\begin{corollary}[Raw finite-record absolute continuity]\label{cor:raw-density}
For every $R\in I$ and $n\ge1$, the actual induced record
$(S_n\kappa_R,S_n r_R,S_n\tau_R^Y)$ under $\nu_R^Y$ has a density with respect to counting measure on $\Z^3$ times Lebesgue measure. For fixed total physical count $N$ and displacement $k$, it has the coarea representation, for almost every $t$,
\begin{equation}\label{eq:raw-coarea}
 p_{n,R}(k,N,t)=\frac1{4\pi\nu(Y_R)}
 \sum_{\mathbf d}\int_{D_{n,\mathbf d}\cap\{L_N=t\}}
              \frac{\dd\mathcal H^1}{|\nabla_{\alpha,p}L_N|}.
\end{equation}
Here $D_{n,\mathbf d}$ is the actual first-return event for the ordered center word $\mathbf d$, total count $N$ and displacement $k$. Bounded measurable insertions have the same formula with the insertion in the numerator. Neither a uniform density bound nor a long-time local limit is asserted by this statement.
\end{corollary}
\begin{proof}
Partition the recurrent regular initial states by total collision count and ordered center word. Uniform finite horizon makes the available lattice steps finite; for fixed $N$ there are finitely many words. On every regular branch the physical length is smooth, and Proposition~\ref{prop:general-edge} makes its critical set empty or a singleton. Away from that null set, the coarea formula applied on a countable exhaustion where $|\nabla L_N|$ is bounded below proves absolute continuity and \eqref{eq:raw-coarea}. Restriction to the measurable first-return event does not alter absolute continuity of the initial measure. The singular initial states have zero flux measure, and summing over the countable possible $N$ gives the full law. The same exhaustion with a bounded insertion proves its signed version.
\end{proof}

For the induced system \eqref{eq:Y}, fix $n$ and restrict its joint law to $\kappa=0$, $r=2n$. Every physical flight is at least $g$. Equality of total length with $2ng$ requires every flight to have length $g$, and hence a nearest-neighbor normal alternating orbit. Of the six possible normal starting states only $(0,0)$ lies in $Y_R$. A neighborhood of this point makes exactly $n$ returns of physical length two each. Compactness, or directly the equality cases of the distance bound, excludes other itineraries from a sufficiently small right neighborhood of $2ng$. Consequently this lattice component of the induced law has right density
\begin{equation}\label{eq:inducededge}
                \frac{J_{2n,R}}{\nu(Y_R)}>0
\end{equation}
at its left endpoint, and zero density on its left.

\begin{corollary}[The boundary term must be retained]\label{cor:nonLone}
For the specified inducing set, the raw joint characteristic function $\Phi_{n,R}$ is not in $L^1(\T^3\times\R)$, for any $n\ge1$. In particular a global integrable bound for this raw function cannot be justified merely by making the normal itinerary exponentially rare.
\end{corollary}
\begin{proof}
If it were integrable, taking its $(0,0,2n)$ lattice Fourier coefficient would give an integrable function of the roof frequency. Its inverse would be a continuous density of that component. A continuous representative cannot agree almost everywhere with zero to the left of $2ng$ and with a function having the positive right limit \eqref{eq:inducededge}. This contradicts Theorem~\ref{thm:edge}.
\end{proof}

The central local-limit problem is not decided by this corollary. The endpoint may be far outside the central window, and $J_{2n,R}$ decays exponentially. What changes is the global inversion method, not the raw observable or the intended central theorem.

\begin{proposition}[Exact edge subtraction]\label{prop:subtract}
Choose a smooth roof cutoff supported in the right neighborhood of Theorem~\ref{thm:edge}, equal to one near its endpoint, and denote the resulting subdensity by $f$. With $J=J_{N,R}$, one has the exact identity
\begin{equation}\label{eq:subtract}
 f(Ng+s)=J e^{-s}\one_{\{s\ge0\}}+q(s),\qquad
 \widehat f(b)=e^{ibNg}\left(\frac{J}{1-ib}+\widehat q(b)\right),
\end{equation}
where $q\in L^1(\R)$, its distributional second derivative is a finite signed measure, and
$|\widehat q(b)|\le C_{N,R}(1+b^2)^{-1}$.
\end{proposition}
\begin{proof}
Writing $f(Ng+s)=\one_{s\ge0}\eta(s)a(s)$, the right value of $\eta a-Je^{-s}$ at zero is zero. Thus $q$ is continuous across zero, piecewise smooth, and exponentially decaying at infinity. Its first derivative may jump at zero, so its second distributional derivative has one finite atom and an integrable density. Distributional integration by parts gives $b^2|\widehat q(b)|\le\|D^2q\|_{\TV}$; the $L^1$ norm gives the bound near zero. Direct integration of the exponential proves \eqref{eq:subtract}.
\end{proof}
This is subtraction and exact addition of an explicit density, not convolution with noise or removal of a physical event. It proves an integrable remainder for the isolated edge. Bounds for all remaining branches and edges are still needed for a long-time result.


\section{A raw residual-inversion theorem}
Let $\eta_{n,R}$ be finite measures on $\Z^3\times\R$ with characteristic functions $\Phi_{n,R}$. Frequencies are $\omega=(u,s,b)\in[-\pi,\pi]^3\times\R$, and our transform convention is $e^{i\omega\cdot z}$. Suppose there are explicit signed edge measures $E_{n,R}$ with densities $e_{n,R}$, and put $H_{n,R}=\Phi_{n,R}-\widehat E_{n,R}$. The following hypotheses are an interface, not consequences of the mechanical constructions above.

\begin{theorem}[Residual criterion for the mixed density LLT]\label{thm:LLT}
Assume the following uniformly for $R\in I$.

\smallskip\noindent
(i) There are $c,C,\delta>0$, $\rho<1$, means $m_R\in\R^4$, and symmetric matrices $cI\le\Sigma_R\le CI$ such that, for $|\omega|\le\delta$,
\[
 \Phi_{n,R}(\omega)=e^{inm_R\cdot\omega}
 \{A_R(\omega)\lambda_R(\omega)^n+O(\rho^n)\},
\]
where $A_R(0)=1$, $|A_R(\omega)-1|\le C|\omega|$,
$\log\lambda_R(\omega)=-\tfrac12\omega^T\Sigma_R\omega+O(|\omega|^3)$, and $|\lambda_R(\omega)|\le e^{-c|\omega|^2}$.

\smallskip\noindent
(ii) $H_{n,R}\in L^1(\T^3\times\R)$ and
\begin{equation}\label{eq:tailcriterion}
 \int_{|\omega|>n^{-2/5}}|H_{n,R}(\omega)|\dd\omega=o(n^{-2}).
\end{equation}

\smallskip\noindent
(iii) $\|E_{n,R}\|_{\TV}=o(n^{-2})$ and, for each fixed $K$,
\[
 n^2\mathop{\rm ess\,sup}_{|(z-nm_R)/\sqrt n|\le K}|e_{n,R}(z)|\longrightarrow0.
\]
Then $\eta_{n,R}$ has a density with respect to counting measure times Lebesgue measure, and on each such central window
\begin{equation}\label{eq:LLT}
 n^2 p_{n,R}(k,m,t)=
 \frac{\exp(-\xi^T\Sigma_R^{-1}\xi/2)}{(2\pi)^2\sqrt{\det\Sigma_R}}+o(1),
 \quad \xi=\frac{(k_1,k_2,m,t)-nm_R}{\sqrt n},
\end{equation}
with uniform essential-supremum error. No smoothing of $\eta_{n,R}$ is used.
\end{theorem}
\begin{proof}
Fourier inversion of $H_{n,R}$ gives a continuous density of the signed measure $\eta_{n,R}-E_{n,R}$. To justify this without assuming a density, take each lattice Fourier coefficient of $H$. It is an integrable roof transform of the corresponding signed component. Convolve that component with a Gaussian approximate identity, use integrability to pass to the inverse-transform limit, and test against compactly supported continuous functions. Uniqueness of finite Fourier transforms identifies each component, hence the full measure. This is a proof of inversion, not a change to the datum. Adding the known density of $E_{n,R}$ gives $p_{n,R}$.

On $|\omega|\le n^{-2/5}$, replacing $H$ by $\Phi$ costs at most a fixed volume times $\|E\|_{\TV}=o(n^{-2})$. Inserting (i), let $v=\sqrt n\,\omega$. The cubic error is uniformly $O(n^{-1/5})$ on this region; the bound on $\lambda$ supplies an integrable Gaussian dominator. Uniform Taylor bounds, truncation first to $|v|\le M$, and then $M\to\infty$, give the inverse Gaussian with error $o(n^{-2})$. The inverse normalization is $(2\pi)^{-4}$ and the Gaussian integral is $(2\pi)^2/\sqrt{\det\Sigma_R}$, giving the constant in \eqref{eq:LLT}. Hypothesis (ii) bounds the remaining frequencies, and (iii) bounds the added density in the central window. All estimates are uniform in $\xi$ in that window.
\end{proof}

For weighted inserted measures, the same proof applies with $A_R(0)=a_R$ and uniformly bounded amplitudes, provided the edge and tail hypotheses are checked for those measures as well. It gives $a_R$ times the Gaussian. It does not assert uniformity for arbitrary $n$-dependent path functionals without such estimates.

\begin{lemma}[An explicit frequency splice]\label{lem:splice}
Suppose a residual transform has a compact minor-arc bound $Ce^{-cn}$ and, for $1\le|b|\le e^{\kappa n}$, a bound $C(1+|b|)^A e^{-cn}$. Beyond $e^{\kappa n}$ suppose it is bounded by
\[
 C_M\rho^n(1+|b|)^{-2}+C_Me^{c_M n}(1+|b|)^{-M}.
\]
If one fixed integer $M>1$ and one $\kappa>0$ satisfy
\begin{equation}\label{eq:splice}
       \frac{c_M}{M-1}<\kappa<\frac{c}{A+1},
\end{equation}
these regions contribute exponentially small $L^1$ mass. Together with the Gaussian major-arc bound between $n^{-2/5}$ and $\delta$, this proves \eqref{eq:tailcriterion}.
\end{lemma}
\begin{proof}
The middle integral is bounded by $C\exp(-(c-(A+1)\kappa)n)$. The two outer integrals are bounded by $C_M\exp((\log\rho-\kappa)n)$ and $C_M\exp(-( (M-1)\kappa-c_M)n)/(M-1)$. The torus has finite volume. The major-arc Gaussian tail is bounded by a polynomial factor times $e^{-c n^{1/5}}$, and the compact remainder is exponential. Strict positivity of the three exponent margins proves the claim.
\end{proof}
The inequality \eqref{eq:splice} must be verified with the actual derivative-growth constants. The phrases ``take $\kappa$ small'' and ``then take $M$ large'' do not imply it: for example, $c_M=2(M-1)$ and $c/(A+1)=1$ leave no possible $\kappa$. This checks a proposed estimate, not the truth or falsity of a central LLT. The residual formulation also requires an independently proved bound for the sum of all extracted edge terms; Proposition~\ref{prop:subtract} treats one physical family only.


\section{Local inversion with nonnegligible remote edge mass}
\label{sec:localized-inversion}
Theorem~\ref{thm:LLT} gives a sufficient global smallness condition on
extracted edges. A central local limit needs less: the extracted measure
must have negligible density and negligible low-frequency convolution
on the central window. Its total variation need not tend to zero.
The following refinement retains the raw datum and the density norm.

Fix an even real function $\chi\in C_c^\infty(\R^4)$ equal to one on
the unit ball and supported in the ball of radius two. Let $a_n\to\infty$
with $a_n=o(n^{1/6})$, put $B_n=a_n/\sqrt n$, and set
$\chi_n(\omega)=\chi(\omega/B_n)$. For all sufficiently large $n$ its
support fits in the fundamental torus interval in the first three
coordinates. With the transform convention already fixed, its inverse
kernel on $\Z^3\times\R$ is
\begin{equation}\label{eq:localized-kernel}
 K_n(k,t)=B_n^4\check\chi(B_nk,B_nt),\qquad
 |K_n(k,t)|\le C_M B_n^4(1+B_n|(k,t)|)^{-M}.
\end{equation}
The check denotes the inverse transform on $\R^4$, including its
$(2\pi)^{-4}$ factor. Convolution here and below uses counting measure
in the discrete coordinates and Lebesgue measure in the roof coordinate.
For fixed $K$, let $W_{n,R,K}=\{z:|z-nm_R|\le K\sqrt n\}$.

\begin{theorem}[A local edge criterion]\label{thm:local-edge-LLT}
Assume the small-frequency hypothesis (i) of Theorem~\ref{thm:LLT}.
Let $E_{n,R}$ be finite signed measures with densities $e_{n,R}$ and
suppose $H_{n,R}=\Phi_{n,R}-\widehat E_{n,R}$ belongs to
$L^1(\T^3\times\R)$. Assume, uniformly in $R$, that
\begin{equation}\label{eq:local-tail}
 \int |1-\chi_n(\omega)|\,|H_{n,R}(\omega)|\dd\omega=o(n^{-2}).
\end{equation}
For every fixed central window suppose also that
\begin{equation}\label{eq:local-edge-errors}
 n^2\mathop{\rm ess\,sup}_{W_{n,R,K}}|e_{n,R}|\longrightarrow0,
 \qquad
 n^2\sup_{W_{n,R,K}}|K_n*E_{n,R}|\longrightarrow0.
\end{equation}
Then the raw density conclusion \eqref{eq:LLT} holds on that window.
There is no small-total-variation hypothesis on $E_{n,R}$.
\end{theorem}
\begin{proof}
As in the proof of Theorem~\ref{thm:LLT}, integrability of $H$ first
identifies a continuous density for $\eta-E$, and adding $e$ gives
an actual mixed density for $\eta$. Splitting this identity with
$\chi_n$ gives the exact equality almost everywhere
\begin{equation}\label{eq:exact-local-decomposition}
 p_{n,R}=K_n*\eta_{n,R}
       +(e_{n,R}-K_n*E_{n,R})
       +\mathcal F^{-1}[(1-\chi_n)H_{n,R}].
\end{equation}
All terms are defined: $E$ is finite, $K_n$ is bounded and integrable,
and the last inverse transform is absolutely convergent.
The last term is at most $(2\pi)^{-4}$ times the integral in
\eqref{eq:local-tail}. The middle term is controlled by
\eqref{eq:local-edge-errors}.

For the first term, its Fourier integral has the factor $\chi_n$.
Insert the assumed small-frequency expansion and use
$v=\sqrt n\,\omega$. The support has $|v|\le2a_n$ and
$nO(|\omega|^3)=O(a_n^3/\sqrt n)=o(1)$. The uniform quadratic
bound on $\lambda_R$ gives a Gaussian dominator, while the amplitude
converges to one. The exponentially small spectral remainder, after
integration and multiplication by $n^2$, is bounded by
$Ca_n^4\rho^n$ and tends to zero. Truncate first at a fixed $|v|$,
then let that truncation grow. This proves uniformly in $R$ and the
central variable that $n^2K_n*\eta_{n,R}$ converges to the Gaussian
in \eqref{eq:LLT}. Combining the three terms proves the assertion.

The kernel is only an analytical device in the exact decomposition
\eqref{eq:exact-local-decomposition}. The two correction terms are
retained and estimated; the result is not a local limit for a
smoothed replacement of the observed law.
\end{proof}

\begin{corollary}[Remote mass need not be small]\label{cor:remote-edges}
Suppose $\|E_{n,R}\|_{\TV}\le C$ and, for the window under
consideration, the support of $E_{n,R}$ is at distance at least
$\varepsilon\sqrt n$ from $W_{n,R,K}$, for a fixed
$\varepsilon>0$. Then \eqref{eq:local-edge-errors} holds.
In particular an extracted density of bounded, nonvanishing mass
can satisfy the local edge hypotheses.
\end{corollary}
\begin{proof}
The density vanishes almost everywhere in the central window. For
$z$ in that window, \eqref{eq:localized-kernel} gives
\[
 n^2|(K_n*E_{n,R})(z)|
 \le CC_M a_n^4(1+\varepsilon a_n)^{-M}\longrightarrow0
\]
by taking $M>4$. The Schwartz bound is available for every fixed $M$
because $\chi$ is smooth. These constants are kernel constants, not
the derivative-growth constants of the billiard in Lemma~\ref{lem:splice}.
\end{proof}

The same conclusion holds for a sum of remote extracted terms whose
total variations are summable with bounded total, and a remaining
nearby extracted density satisfying \eqref{eq:local-edge-errors}.
This follows by linearity and the triangle inequality. Thus individual
edge coefficients do not all have to be $o(n^{-2})$ merely because
one proves a central limit. Nevertheless, one must still construct
the full extraction and prove \eqref{eq:local-tail}. Neither this
corollary nor the single-word estimate in Proposition~\ref{prop:general-edge}
controls the sum of all central critical words or singular boundaries.
A restriction of a density to a sharp window can itself create new
jump terms; it is not automatically an integrable residual.

For weighted measures the same proof gives the corresponding amplitude
times the Gaussian, provided the small-frequency expansion, residual
estimate and local edge errors are verified for that weighted class.
This is the precise additional input needed before applying
Proposition~\ref{prop:conditioning} to a proposed family of insertions.

\section{Conditioning, clocks, and geometric prediction}
\begin{proposition}[Central density conditioning]\label{prop:conditioning}
Suppose an unweighted and an $f$-weighted joint density satisfy
\[
 n^2p_n(z)=G_n(z)+r_n(z),\qquad
 n^2p_n^f(z)=a_fG_n(z)+r_n^f(z)
\]
on a common central window, with $G_n\ge g_0>0$ and $|r_n|,|r_n^f|\le\varepsilon_n<g_0$. Then their density ratio differs from $a_f$ by at most
\begin{equation}\label{eq:ratio}
          \frac{(1+|a_f|)\varepsilon_n}{g_0-\varepsilon_n}.
\end{equation}
The same bound holds after conditioning on fixed lattice coordinates and any positive-length roof interval contained in that window, including intervals whose lengths shrink with $n$.
\end{proposition}
\begin{proof}
Subtract $a_f$ times the first identity from the second and divide by $G_n+r_n\ge g_0-\varepsilon_n$. For an interval, integrate both identities first; the errors are at most its length times $\varepsilon_n$ and the denominator is at least its length times $g_0-\varepsilon_n$. The length cancels.
\end{proof}
At an exact roof value this is a density-defined version of a regular conditional expectation for almost every value with positive density. It is not conditioning on a positive-probability singleton. It supplies no estimate outside the common central window and no functional path convergence by itself. These are the precise limitations of this interface to stopped-path large deviations.

\begin{proposition}[A conditional physical-clock transfer]\label{prop:clock}
Let $(X_j,\tau_j)\in\R^d\times(0,\infty)$ be stationary dynamical observations with means $(a,\mu)$, $\mu>0$, and a joint functional CLT with covariance $\Sigma$. Assume the uniform law of large numbers for their clock and that initial and terminal residual records on bounded rescaled time intervals are $o_P(\sqrt t)$. Then the reward accumulated up to physical time $t$, centered by $(a/\mu)t$, has limiting covariance per unit time
\begin{equation}\label{eq:clock}
 \frac1\mu Q\Sigma Q^T,\qquad Q=(I_d,-a/\mu).
\end{equation}
For $X=(\kappa_1,\kappa_2,r)$ this transfers induced displacement and collision counts, not the number of experimental bits.
\end{proposition}
\begin{proof}
The clock law gives its inverse index $N(t)/t\to1/\mu$. The centered functional CLT evaluated at this inverse time and the identity $S_{N(t)}\tau=t+o_P(\sqrt t)$ give
\[
 S_{N(t)}X-(a/\mu)t
 = (S_{N(t)}X-N(t)a)-(a/\mu)(S_{N(t)}\tau-N(t)\mu)+o_P(\sqrt t).
\]
The continuous limit permits the random-time substitution. Projection by $Q$ and evaluation at time $1/\mu$ prove \eqref{eq:clock}. The assumed residual estimate includes unfinished induced returns, which are not bounded merely because individual flights are bounded.
\end{proof}

\begin{lemma}[Covariance continuity from a finite-time comparison]\label{lem:covariance}
Suppose a centered correlation matrix satisfies $\|C_j(R)\|\le Ce^{-cj}$ and
$\|C_j(R)-C_j(R')\|\le C\delta^\alpha e^{a j}$, where $\delta=|R-R'|$, $a,c,\alpha>0$. Then the absolutely convergent Green--Kubo covariance satisfies
\[
 \|\Sigma_R-\Sigma_{R'}\|\le C'\delta^{\alpha c/(a+c)}.
\]
If $a=0$, the bound is instead $C'\delta^\alpha(1+|\log\delta|)$.
\end{lemma}
\begin{proof}
Split the defining correlation series at $J=\lfloor\alpha\log(1/\delta)/(a+c)\rfloor$. For $a>0$, the finite sum is bounded by $C\delta^\alpha e^{aJ}$ and the remaining sum by $Ce^{-cJ}$. Both have the stated order. For $a=0$, sum the constant comparison bound over $J+1$ terms and use the exponential tail. Symmetrizing the correlations only changes the constant.
\end{proof}
The hypotheses of this lemma must concern the actual induced observables and their varying domains. Spectral continuity for a different observable norm is not a substitute for this comparison.

\begin{proposition}[Geometric error and the central prediction scale]\label{prop:prediction}
Assume \eqref{eq:LLT} with error $\epsilon_n$, uniformly elliptic covariances, and estimates $(\widehat m,\widehat\Sigma)$ satisfying $|\widehat m-m_R|+\|\widehat\Sigma-\Sigma_R\|\le\omega(\nu)$. If $\sqrt n\,\omega(\nu)\le1$, replacing the Gaussian parameters in \eqref{eq:LLT} changes its right side on a fixed central window by at most $C(1+\sqrt n)\omega(\nu)$. The total prediction error is bounded by this quantity plus $\epsilon_n$.
\end{proposition}
\begin{proof}
The argument of the Gaussian moves by $\sqrt n(\widehat m-m_R)$. On a slightly enlarged compact window the Gaussian and its derivatives in its argument and covariance are bounded uniformly under ellipticity. The mean value theorem gives the bound. Adding the LLT error finishes the proof.
\end{proof}

There is also a direct equilibrium conclusion that needs no local-limit hypothesis. The mean collision rate is
\[
 \lambda(R)=\frac{2R}{\sqrt3/2-\pi R^2}.
\]
Its derivative is $2/A_R+4\pi R^2/A_R^2$, where $A_R=\sqrt3/2-\pi R^2$. On $I$, $A_R>17/100$, using $\sqrt3>173/100$ and $\pi<22/7$, and direct substitution gives $0<\lambda'(R)<110$. Therefore a radius estimate with error at most $\nu$ gives collision-rate error at most $110\nu$. On the circular submodel the average support function identifies the radius, independently of translation. The finite geometric conclusion of \cite{A2G} therefore transfers to this rate with its geometric observation cost, without recovering the entire unknown launch law.


\section{Exact return means and stationary unfinished flights}
\label{sec:exact-return-means}
The enlarged section gives explicit normalization constants without any
local-limit assumption. We use ergodicity of the finite-horizon
periodic dispersing collision map, as supplied by the standard billiard
framework \cite{Young}. No spectral conclusion for the new first-return
map is imported from that fact.

\begin{proposition}[Exact induced means]\label{prop:exact-induced-means}
For the actual return records to $Y_R^*$,
\begin{equation}\label{eq:induced-mean-vector}
 \int(\kappa_{R,1}^*,\kappa_{R,2}^*,r_R^*,\tau_R^*)\dd\nu_R^*
 =\left(0,0,\frac1{c_*},\frac{\bar\tau_R}{c_*}\right),
 \quad
 \bar\tau_R=\frac{\sqrt3}{4R}-\frac{\pi R}{2}.
\end{equation}
The mean roof satisfies
\begin{equation}\label{eq:mean-roof-modulus}
 \left|\frac{\mathrm d}{\mathrm dR}\int\tau_R^*\dd\nu_R^*\right|
 <\frac{15}{4c_*}\qquad(R\in I).
\end{equation}
The equilibrium physical collision rate remains $1/\bar\tau_R$.
\end{proposition}
\begin{proof}
For any integrable physical collision observable $f$, the disjoint
return-tower levels partition the collision section up to null sets.
Invariance, first for nonnegative functions and then for positive and
negative parts, gives the Kac identity
\[
 \int_{Y_R^*}\sum_{j=0}^{r_R^*(x)-1} f(T_R^j x)\dd\nu(x)
                         =\int f\dd\nu.
\]
Taking $f=1$ and dividing by $c_*$ gives the mean return count.
Taking $f=\tau_R$ gives the mean induced roof. The physical
suspension normalization proved earlier gives the displayed value
of $\bar\tau_R$. The physical displacement is integrable by the
finite-horizon bound. Spatial inversion of the entire lattice maps
its value to its negative and preserves $\nu$, so its mean is zero.
This proves \eqref{eq:induced-mean-vector}.

The section volume $c_*$ is independent of $R$. Differentiating the
explicit expression for $\bar\tau_R$ yields
$-\sqrt3/(4R^2)-\pi/2$. The bounds $\sqrt3<7/4$,
$\pi<22/7$ and $R\ge9/20$ show its absolute value is less than
$15/4$. Finally the ratio of mean physical collisions per return to
mean physical time per return is $1/\bar\tau_R$.
\end{proof}

\begin{proposition}[The complete induced suspension]\label{prop:induced-suspension}
The equilibrium flow is represented by the suspension over $F_R^*$
with roof $\tau_R^*$ and probability
\begin{equation}\label{eq:induced-suspension}
 \frac{\nu_R^*(\mathrm dx)\,\mathrm ds}{\int\tau_R^*\dd\nu_R^*},
               \qquad 0\le s<\tau_R^*(x).
\end{equation}
If $A_R$ and $B_R$ denote respectively the age since the last visit
and the time to the next visit to $Y_R^*$ under this stationary law,
then, for $q\ge0$,
\begin{equation}\label{eq:residual-tails}
 \Pr(A_R>q)=\Pr(B_R>q)
 =\frac{\int(\tau_R^*-q)_+\dd\nu_R^*}
              {\int\tau_R^*\dd\nu_R^*}.
\end{equation}
At either one fixed endpoint of a time interval of length $t$,
the incomplete induced flight, its physical collision count and its
lattice displacement divided by $\sqrt t$ converge to zero in
probability as $t\to\infty$, for each fixed $R$.
\end{proposition}
\begin{proof}
In the physical suspension over $T_R$, group the consecutive
single-flight intervals before each first return to $Y_R^*$.
The tower identity and the normalization
$c_*\int\tau_R^*\dd\nu_R^*=\bar\tau_R$ convert the original
suspension measure exactly to \eqref{eq:induced-suspension}.
This rearranges the actual dependent path; it introduces no new
launch law and no independent flight variables. The marked formula
of Proposition~\ref{prop:marked} can consequently be grouped into
these complete return records while keeping both endpoint pieces.

For a given roof $\tau$, the set of $s\in[0,\tau)$ with age $s>q$
has length $(\tau-q)_+$. The set with remaining time $\tau-s>q$
has the same length. Integration proves \eqref{eq:residual-tails}.
Its right side tends to zero by the integrability in
Proposition~\ref{prop:exact-induced-means}. Each endpoint roof is
finite almost surely under the length-biased stationary law;
stationarity gives the same marginal at the other fixed endpoint.
Thus its ratio to $\sqrt t$ tends to zero in probability.

Every physical flight has length at least $1-2R\ge3/50$. The
number of collisions in an incomplete interval of length $s$ is
at most $1+s/(1-2R)$. Its displacement between fundamental cells
is bounded by a constant plus a fixed multiple of $s$, since the
speed is one and the lattice cell has bounded diameter. These
bounds prove the same endpoint statement for the two marked rewards.
\end{proof}

A fixed-endpoint assertion is not a bound on the maximum residual over
all times in $[0,t]$. Nor does pointwise integrability in $R$ imply a
uniform residual estimate over $I$. Such uniform or process-level
claims require uniform integrability or a quantitative tail for the
induced roof. The functional limit assumptions in
Proposition~\ref{prop:clock} therefore remain explicit. The exact
mean and the modulus \eqref{eq:mean-roof-modulus}, by contrast, are
already available for geometric error propagation. They do not
require recovery of the entire preparation distribution in A2-GEOM.

\section{Stability of the actual return records}
\label{sec:return-stability}
We compare the first returns themselves, including their changing domains.
The comparison is made on the common collision space
$M=\T\times(-1,1)$ with the probability $\nu$ in \eqref{eq:flux}.
It is not a comparison between independently chosen inducing schemes.
The squares defining $Y_R$ move analytically with $R$, and $U$ is fixed.
Consequently $Y_R^*$ has constant measure $c_*$, and its boundary is a
finite union of smooth arcs, continuously varying with $R$.

We use only ergodicity, not a parameter-uniform mixing rate, from the
finite-horizon Sinai theory. In the present family the obstacles are
smooth, strictly dispersing and disjoint, and the horizon is finite.
The flow-mixing theorem \cite[Corollary 1.3 and correction note]{BDL}
therefore implies ergodicity of the collision map: the suspension of a
nontrivial invariant collision set would be a nontrivial invariant flow
set. Thus every $Y_R^*$ has full saturation. This explicitly supplies the
full-tower hypothesis in the Kac identity used above.

Write $r=r_R^*$, $F=F_R^*$ and $\tau^*=\tau_R^*$ in sums indexed by
induced returns. For $n\ge1$ define the actual cumulative record
\[
 J_{n,R}=\left(\sum_{j<n}\kappa_R^*\circ F^j,
                   \sum_{j<n}r_R^*\circ F^j,
                   \sum_{j<n}\tau_R^*\circ F^j\right)
       =(K_{n,R},N_{n,R},T_{n,R})
\]
on $Y_R^*$. Its values lie in $\Z^3\times\R$; in particular the last
coordinate is one cumulative time, not $n$ independent roof coordinates.
A tilde denotes extension by zero to $M$. Let $p_{n,R}$ denote the raw
mixed density of $J_{n,R}$ under $\nu_R^*$, whose existence is included
in Theorem~\ref{thm:return-stability} below. Norms of densities use
counting measure times Lebesgue measure, denoted by $\mathfrak m$.

\begin{lemma}[Almost-everywhere stability of a finite return record]
\label{lem:finite-return-stability}
Fix $R_0\in I$ and a finite $n$. Outside a $\nu$-null subset of $M$,
$\widetilde J_{n,R}(x)\to\widetilde J_{n,R_0}(x)$ as $R\to R_0$.
For $x\in Y_{R_0}^*$ outside that null set, the physical word, every
intermediate return decision, $K_{n,R}(x)$ and $N_{n,R}(x)$ are constant
for all sufficiently close $R$. On a neighborhood of such a point the
cumulative time is the physical analytic length of that fixed word.
The neighborhood and parameter tolerance may depend on $x$ and $n$.
\end{lemma}
\begin{proof}
Remove the boundary of $Y_{R_0}^*$, all singular initial states and all
preimages under $T_{R_0}$ of that boundary. These sets are null: the
boundary has zero area, the map preserves $\nu$, and the billiard
singularities and their iterates have zero flux measure. Remove also
the null set on which a required return is infinite. For a remaining
$x\in Y_{R_0}^*$, the first $n$ returns contain finitely many regular
physical collisions. Every tested state lies strictly inside or strictly
outside the return set. The finite-itinerary analyticity following
\eqref{eq:entry}, and the continuous motion of the return boundary,
preserve all these decisions for nearby $(R,x)$. They also preserve
the selected entry roots and lattice labels. The final physical count
and displacement are therefore unchanged, while the sum of flight
lengths varies analytically. A remaining point outside $Y_{R_0}^*$
stays outside $Y_R^*$ for nearby $R$, and its zero extension stays zero.
These statements prove the asserted convergence. No common neighborhood
for infinitely many words is used.
\end{proof}

\begin{theorem}[Finite-record stability in density and first moment]
\label{thm:return-stability}
For every fixed $n\ge1$, the maps
\begin{equation}\label{eq:return-Lone}
 R\longmapsto\widetilde J_{n,R}\in L^1(\nu;\R^4),\qquad
 R\longmapsto p_{n,R}\in L^1(\mathfrak m)
\end{equation}
are continuous. Both maps are uniformly continuous on $I$.
Moreover $\{\widetilde N_{n,R}:R\in I\}$ and
$\{\widetilde T_{n,R}:R\in I\}$ are uniformly integrable.
If $\eta_{n,R}$ is the record law, then
\begin{equation}\label{eq:record-Wone}
 W_1(\eta_{n,R},\eta_{n,S})
 \le c_*^{-1}\|\widetilde J_{n,R}-\widetilde J_{n,S}\|_{L^1(\nu)},
\end{equation}
where $W_1$ uses the Euclidean distance on $\Z^3\times\R$.
In particular, for the original, unmodified characteristic functions,
\begin{align}\label{eq:raw-frequency-comparison}
 |\Phi_{n,R}(\omega)-\Phi_{n,S}(\omega)|
 &\le \|p_{n,R}-p_{n,S}\|_{L^1(\mathfrak m)},\nonumber\\
 |\Phi_{n,R}(\omega)-\Phi_{n,S}(\omega)|
 &\le c_*^{-1}|\omega|\,
       \|\widetilde J_{n,R}-\widetilde J_{n,S}\|_{L^1(\nu)}.
\end{align}
No modulus uniform in $n$, and no frequency-integrated estimate, is
asserted by these inequalities.
\end{theorem}
\begin{proof}
The Kac identity and invariance of $F_R^*$ give, exactly,
\begin{equation}\label{eq:block-Kac-mass}
 \int_M\widetilde N_{n,R}\dd\nu=n,\qquad
 \int_M\widetilde T_{n,R}\dd\nu=n\bar\tau_R.
\end{equation}
For nonnegative functions $h_j\to h$ almost everywhere with
$\int h_j\to\int h$, the identity
\[
 \|h_j-h\|_1=\int h_j+\int h-2\int\min(h_j,h)
\]
and dominated convergence for $\min(h_j,h)\le h$ prove $L^1$
convergence. Apply this to the two nonnegative coordinates, using
Lemma~\ref{lem:finite-return-stability}, \eqref{eq:block-Kac-mass}
and continuity of $\bar\tau_R$. Uniform finite horizon bounds every
single-flight lattice label by a fixed constant. Thus
$|\widetilde K_{n,R}|\le C\widetilde N_{n,R}$. The already established
$L^1$ convergence of the counts makes these bounds uniformly integrable
along each convergent parameter sequence. Almost-everywhere convergence
and truncation then give $L^1$ convergence of the displacement as well.
This proves the first continuity in \eqref{eq:return-Lone}.

A continuous image of the compact interval $I$ is compact in $L^1$.
For completeness, an $L^1$-compact nonnegative family is uniformly
integrable. Cover it by finitely many $L^1$ balls of radius $\epsilon$
with centers $h_1,\ldots,h_q$. For a member $h$ in the ball about $h_i$,
\[
 \int_{\{h>A\}}h
 \le 2\|h-h_i\|_1+
              2\int_{\{h_i>A/2\}}h_i.
\]
Choose $A$ for the finite set of centers and then let $\epsilon$
decrease. This proves the stated uniform integrability without a tail
assumption or an exponential return estimate.

We next prove existence and $L^1$ continuity of the density. Partition
$Y_R^*$, up to null sets, by total physical count and ordered lattice
word. The argument of Corollary~\ref{cor:raw-density} applies to this
actual first-return event as well: on each regular word the cumulative
length has at most one critical point by
Proposition~\ref{prop:general-edge}; elsewhere it is a submersion.
Coarea on a countable exhaustion and then summation give absolute
continuity. This argument recomputes the return event for $Y_R^*$;
it does not transfer a density normalization from $Y_R$.

Fix $R_0$ and $\epsilon>0$. After deleting the null sets just described,
choose finitely many smooth nonnegative cutoffs $\psi_i$ on $M$, with
$\sum_i\psi_i\le1$, whose supports lie compactly in regular
first-$n$-return patches for $R_0$, such that
\[
 c_*^{-1}\int\sum_i\psi_i\dd\nu>1-\epsilon.
\]
Refine the patches so that on each one a coordinate derivative of the
cumulative length is nonzero. Compact support and the finite-return
lemma preserve the same word, all return decisions and this derivative
for $R$ close to $R_0$. On each patch the map consisting of the other
initial coordinate and the cumulative length is a local diffeomorphism.
It and its inverse depend continuously in $C^1$ on $R$ on slightly
larger compact patches. Changing variables and integrating the other
coordinate therefore shows that the pushforward of
$c_*^{-1}\psi_i\nu$ has a density continuous in $L^1$ with $R$.
One can see this directly by extending the transformed smooth weight
by zero off its coordinate patch: the compactly supported cutoff
vanishes near the patch boundary, the Jacobian is bounded away from
zero, and dominated convergence applies on a fixed bounded coordinate
box. The lattice coordinates are constant on the patch.

The remainder is a positive measure. Its mass is
$1-c_*^{-1}\int\sum_i\psi_i\dd\nu<\epsilon$, for all these $R$,
because the cutoff supports stay inside $Y_R^*$ and $\nu(Y_R^*)=c_*$.
Consequently
\[
 \|p_{n,R}-p_{n,R_0}\|_1
 \le\sum_i\|p_{i,R}-p_{i,R_0}\|_1+2\epsilon.
\]
Let $R\to R_0$ and then $\epsilon\to0$. This proves the second
continuity in \eqref{eq:return-Lone}. Compactness of $I$ gives uniform
continuity of both maps.

For a 1-Lipschitz function $f$, subtract the constant $f(0)$. Then
\[
 \int f\dd\eta_{n,R}
 =c_*^{-1}\int_M f(\widetilde J_{n,R})\dd\nu.
\]
The same identity holds for $S$. Subtraction and the Lipschitz bound,
followed by the dual characterization of $W_1$, give
\eqref{eq:record-Wone}. All first moments are finite by
\eqref{eq:block-Kac-mass}. Direct integration against the densities
gives the first bound in \eqref{eq:raw-frequency-comparison}. For the
second, use the zero extensions and
$|e^{ia}-e^{ib}|\le|a-b|$; the off-section constants cancel because
both section measures equal $c_*$. These are estimates of the raw
transforms, with no convolution or change of observable.
\end{proof}

\begin{proposition}[Weighted records and a positive denominator]
\label{prop:weighted-return-stability}
Fix $n$. Suppose measurable insertions $w_R:M\to\R$ satisfy
$|w_R|\le M_0$ and, for every $R_0$, $w_R(x)\to w_{R_0}(x)$ for
$\nu$-almost every $x$ as $R\to R_0$. The $w_R$-weighted record
densities $p^w_{n,R}$ are continuous in $L^1(\mathfrak m)$.
Set
\[
 d_n(\delta)=\sup_{|R-S|\le\delta}\|p_{n,R}-p_{n,S}\|_1,
 \qquad
 d_n^w(\delta)=\sup_{|R-S|\le\delta}\|p^w_{n,R}-p^w_{n,S}\|_1.
\]
These moduli tend to zero. For any Borel observation event $B$ in the
record space, if $\eta_{n,R}(B)\ge b>0$ and
$d_n(\delta)<b$, then, whenever $|R-S|\le\delta$,
\begin{equation}\label{eq:finite-conditioning-stability}
 \left|\frac{\int_Bp^w_{n,R}\dd\mathfrak m}{\eta_{n,R}(B)}
       -\frac{\int_Bp^w_{n,S}\dd\mathfrak m}{\eta_{n,S}(B)}\right|
 \le\frac{d_n^w(\delta)+M_0d_n(\delta)}{b-d_n(\delta)}.
\end{equation}
\end{proposition}
\begin{proof}
Use the compact regular patches in the preceding proof. Replacing
$w_R$ by $w_{R_0}$ costs at most
$c_*^{-1}\int|w_R-w_{R_0}|\dd\nu\to0$, by dominated convergence
and contraction of total variation under pushforward. On each fixed
patch approximate the bounded measurable weight $w_{R_0}$ in $L^1$
by a smooth one. The approximation error in its pushforward is at
most the same $L^1$ error, uniformly in $R$. For the smooth weight,
the coordinate-change proof above applies. The leftover measure has
total variation at most $M_0\epsilon$. This proves weighted density
continuity, and compactness gives the modulus.

Write $a_R=\int_Bp^w_{n,R}\dd\mathfrak m$ and
$b_R=\eta_{n,R}(B)$. Then
$|a_R-a_S|\le d_n^w(\delta)$,
$|b_R-b_S|\le d_n(\delta)$ and $|a_R|\le M_0b_R$.
Since $b_S\ge b-d_n(\delta)>0$, subtraction of the two ratios gives
\eqref{eq:finite-conditioning-stability}.
\end{proof}

The event $B$ may fix the lattice coordinates and restrict the roof
to an interval. The denominator in
\eqref{eq:finite-conditioning-stability} is an actual event probability.
The estimate does not remove its deterioration for small events. In
particular density $L^1$ continuity alone is not the shrinking-interval
LLT of Proposition~\ref{prop:conditioning}; that proposition retains
its stronger pointwise density hypotheses. An error bound
$|\widehat R-R|\le\delta$ can be inserted into the displayed moduli,
but no numerical sampling budget follows from a noneffective modulus.

\begin{corollary}[Uniform stationary endpoint control]
\label{cor:uniform-stationary-endpoints}
Put $H_R=\widetilde T_{1,R}$ and
$b_*=\min_{R\in I}\bar\tau_R>0$. The stationary induced suspensions,
regarded as probabilities on $M\times\R_+$, have densities
\[
 h_R(x,s)=\bar\tau_R^{-1}\one_{\{0\le s<H_R(x)\}}
\]
with respect to $\nu(\dd x)\dd s$, and
\begin{equation}\label{eq:suspension-Lone}
 \|h_R-h_S\|_1\le
 b_*^{-1}\bigl(\|H_R-H_S\|_1+|\bar\tau_R-\bar\tau_S|\bigr).
\end{equation}
Furthermore
\begin{equation}\label{eq:uniform-roof-tail}
 \Psi(q):=b_*^{-1}\sup_{R\in I}
                   \int_M H_R\one_{\{H_R>q\}}\dd\nu
             \longrightarrow0\quad(q\to\infty).
\end{equation}
Under the stationary flow, at either deterministic endpoint of an
interval of length $t$, the incomplete induced time, physical collision
count and lattice displacement are $o_P(\sqrt t)$, uniformly over
$R\in I$. The same is true jointly at any fixed finite number of
deterministic endpoints.
\end{corollary}
\begin{proof}
The normalization in \eqref{eq:induced-suspension} and
$c_*\int\tau_R^*\dd\nu_R^*=\bar\tau_R$ give $h_R$ exactly.
For fixed $x$, the symmetric difference of the intervals
$[0,H_R(x))$ and $[0,H_S(x))$ has length $|H_R(x)-H_S(x)|$.
Integrating, and then comparing the two normalizing constants, proves
\eqref{eq:suspension-Lone}. The uniform integrability in
Theorem~\ref{thm:return-stability} proves
\eqref{eq:uniform-roof-tail}.

The full roof containing a stationary endpoint has tail
\[
 \Pr_R\{\text{containing roof}>q\}
 =\bar\tau_R^{-1}\int_MH_R\one_{\{H_R>q\}}\dd\nu
 \le\Psi(q).
\]
This is the length-biased law; replacing it by
$\nu_R^*(\tau_R^*>q)$ would be incorrect. Both endpoint pieces
are at most this full roof. The collision count of an incomplete
piece of length $s$ is at most $1+(50/3)s$, and its lattice
displacement is bounded by $C(1+s)$, by the physical estimates in
Proposition~\ref{prop:induced-suspension}. Hence there are fixed
constants $C_0,C_1$ such that each marked endpoint norm is bounded
by $C_0+C_1H$ when its containing roof has length $H$. Its probability of exceeding
$\epsilon\sqrt t$ is at most
$\Psi((\epsilon\sqrt t-C_0)/C_1)$ for large $t$, uniformly in $R$.
Stationarity applies at every deterministic endpoint, including an
endpoint depending deterministically on $t$. A union bound proves
the finite-endpoint assertion; independence is not used.
\end{proof}

This removes the additional uniform-integrability premise for fixed
stationary endpoints in this family. It is not a bound on a supremum
over a growing time interval or at an arbitrary stopping time. The
functional clock theorem still needs its printed process hypotheses.
Likewise \eqref{eq:raw-frequency-comparison} is a uniform comparison
of raw characteristic functions, not an integrable tail for any one
of them. The high-frequency and central all-branch remainder problem
is unchanged by an $L^1$ comparison of probability densities.

\section{Exponential tails for the genuine return record}
\label{sec:exponential-returns}
We now use the spectral theory of the collision map, in addition to the
geometric and ergodic inputs used above. The operator in this section is
not the operator of the induced map. A smooth killing weight on the
collision space gives a scalar estimate for avoiding the actual return
section. In particular no multiplication by its discontinuous indicator
on an anisotropic space will be needed.

\subsection{The collision-space input}
We use the strong and weak distribution spaces of Demers--Zhang
\cite[Section 3.3]{DZ}. Their properties needed here are the local uniform
spectral gap and power bounds for unforced dispersing collision maps,
including small changes of the scatterers and their boundary lengths
\cite[Theorems 2.1--2.6 and Remark 2.7(b)]{DZ}. The spectral gap for an
individual unforced Lorentz map is the input from \cite{DZ11} used there.
Here is the precise consequence, together with the checks for this family.

\begin{lemma}[Local uniform collision spaces]
\label{lem:collision-spectral-input}
There is a finite cover of $I$ by parameter neighborhoods $I_j$. On each
$I_j$ there is a Banach space $\mathcal B_j$ of distributions on a common
collision space, a continuous mass functional $\ell_j$, and transfer
operators $\mathcal L_R$, such that
\begin{equation}\label{eq:collision-split}
 \mathcal L_R=\Pi_R+\mathcal N_R,\quad
 \Pi_R h=\ell_j(h)\nu,\quad
 \Pi_R\mathcal N_R=\mathcal N_R\Pi_R=0,\quad
 \|\mathcal N_R^k\|_{\mathcal B_j}\le C\vartheta^k
 \quad(k\ge0),
\end{equation}
where $0<\vartheta<1$ and $C<\infty$ can be chosen for all the finitely
many neighborhoods. The norms of $\nu$, $\ell_j$ and $\mathcal L_R$ are
uniformly bounded. Multiplication by a smooth function $g$ satisfies
\begin{equation}\label{eq:smooth-multiplier}
 \|g h\|_{\mathcal B_j}\le C\|g\|_{C^2}\|h\|_{\mathcal B_j}.
\end{equation}
The transfer convention is $(\mathcal L_R h)(f)=h(f\circ T_R)$.
No assertion about the unbounded induced twists is included in this lemma.
\end{lemma}
\begin{proof}
For the collision-space construction use coordinates $(\alpha,\varphi)$,
where $p=\sin\varphi$. The invariant probability is the same for all
parameters: $\dd\nu=(4\pi)^{-1}\cos\varphi\dd\alpha\dd\varphi$.
We do not conjugate a distribution space by the inverse of
$p=\sin\varphi$ at grazing. At a fixed $R_0$, use the boundary coordinate
$r_0=R_0\alpha$. The parametrization of the nearby circle is then
$r_0\mapsto R n_{r_0/R_0}$. It is a $C^3$-bounded deformation, $C^2$-close
to the original parametrization, and its constant speed $R/R_0$ tends to
one. This is precisely the boundary-length reparametrization permitted in
\cite[Remark 2.7(b)]{DZ}.

The physical gap between different disks is at least $3/50$. The
curvatures belong to $[100/47,20/9]$, and the parametrized boundaries have
uniform $C^3$ bounds. The finite horizon is uniform on $I$: a line with
no disk intersection at the smallest radius would be an infinite free
line in that table, which has no corridor; compactness of position and
direction then bounds the free length, and increasing the radius cannot
increase it. If a rectangular fundamental domain is desired, use the
two-scatterer cover with periods $(1,0)$ and $(0,\sqrt3)$ and disk centers
$0,b$. Its deck translation commutes with the collision map. Restricting
to the deck-invariant distributions gives the original table. Thus the
local billiard estimates are applied without an affine distortion of
specular reflection.

The uniform geometric hypotheses of \cite[Theorem 2.5]{DZ} and the
perturbation hypothesis of its Theorem 2.6 are satisfied. The unforced
collision map at each $R_0$ has the simple eigenvalue one and a spectral
gap. The strong/weak perturbation theorem, specifically the complementary
power bound in \cite[Theorem 2.1(3)]{DZ}, gives \eqref{eq:collision-split}
on a neighborhood of $R_0$. The eigenprojection has the displayed form
because the invariant probability is $\nu$ and mass is preserved. A
finite cover of the compact parameter interval makes the constants
uniform. This uses local common spaces and does not assert that all
arbitrary inducing schemes share a Banach space.

For completeness, smooth multiplication follows directly from the
stable-test and unstable-comparison norms defining these spaces. On a
single admissible stable curve, multiplying its test function by $g$
changes its H\"older norm by at most $C\|g\|_{C^1}$. For two matched
curves at distance $\epsilon$, split the difference of the weighted tests
into the old test difference and the difference of the two restrictions
of $g$. After graph matching, the latter has $C^q$ norm at most
$C\|g\|_{C^2}\epsilon^{1-q}$ by interpolation between its $C^0$ and
$C^1$ bounds. The defining exponents satisfy $\beta<1-q$, so division
by $\epsilon^\beta$ is bounded for $\epsilon\le1$. The stable part of the
norm controls this second term and the unstable part controls the first.
The weak norm has the same single-curve multiplication bound. Density of
the defining smooth distributions extends these bounds to
$\mathcal B_j$, proving \eqref{eq:smooth-multiplier}. The coordinate
changes above and the finite cover have uniformly bounded derivatives.
\end{proof}

\subsection{A smooth avoidance estimate}
Choose once and for all $\chi\in C^\infty(M)$ such that
\begin{equation}\label{eq:fixed-killing-weight}
 0\le\chi\le1,\quad
 \supp\chi\Subset\{ |\alpha|<1/400,\ |p|<1/400\},\quad
 a_\chi:=\int\chi\dd\nu>0.
\end{equation}
The indicated rectangle is inside the square of $Y_R$ centered at
$(0,0)$ for every $R$; hence it is inside $Y_R^*$. Its support is away
from grazing, so it is smooth also in the collision coordinates used in
Lemma~\ref{lem:collision-spectral-input}. This is a proof weight, not an
extra field in the collision record.

\begin{proposition}[Uniform soft-killing estimate]
\label{prop:soft-killing}
There exist $z_0,c,C>0$ such that for every $R\in I$ and $k\ge0$,
\begin{equation}\label{eq:soft-killing}
 \int_M\exp\left(-z_0\sum_{j=0}^{k-1}\chi\circ T_R^j\right)\dd\nu
       \le Ce^{-ck}.
\end{equation}
\end{proposition}
\begin{proof}
On a local collision space set
$\mathcal L_{R,z}=\mathcal L_R M_{e^{-z\chi}}$ for complex $z$.
The multiplier bound proves analyticity in the strong operator norm and,
uniformly in $R$, $\|\mathcal L_{R,z}-\mathcal L_R\|\le C|z|$ near
zero. Equation~\eqref{eq:collision-split} bounds the unperturbed
resolvents uniformly on a circle around one and on a second circle
$|w|=\vartheta_1$ with $\vartheta<\vartheta_1<1$, chosen disjointly.
A uniformly convergent resolvent Neumann series therefore defines a
rank-one analytic projection and a single eigenvalue $\lambda_R(z)$
near one for all $|z|<z_*$. The complementary powers are bounded by
$C\vartheta_1^k$, after decreasing $z_*$ if necessary.

Since $\ell_j(\nu)=1$, differentiation at zero gives
\[
 \lambda_R(0)=1,\qquad
 \lambda_R'(0)=-\ell_j(\mathcal L_R(\chi\nu))=-a_\chi.
\]
The fixed complex disk and the uniform resolvent bounds give a uniform
bound on the second derivative of $\lambda_R$ on a smaller disk, by
Cauchy's formula. Thus
$|\lambda_R(z)-1+a_\chi z|\le C_2|z|^2$ uniformly in $R$.
For real small $z$ the eigenvalue is real, by conjugation and uniqueness,
and positive, by its closeness to one. Choose a fixed $z_0>0$ such that
$C_2z_0\le a_\chi/2$ and $z_0<z_*$. Then
$0<\lambda_R(z_0)\le1-a_\chi z_0/2$. The spectral decomposition gives
$\|\mathcal L_{R,z_0}^k\|\le Ce^{-ck}$. Only finitely many collision
spaces occur, so the same $z_0,c,C$ work on all of $I$.

Iteration of the transfer identity, first for densities and then as
bounded functionals, yields
\[
 \ell_j(\mathcal L_{R,z_0}^k\nu)
 =\int\prod_{j=0}^{k-1}e^{-z_0\chi(T_R^j x)}\dd\nu(x).
\]
The norms of $\ell_j$ and $\nu$ are bounded. This proves the asserted
scalar estimate. In particular the argument does not apply a local
large-deviation formula outside its stated neighborhood of the mean.
\end{proof}

\begin{theorem}[Exponential moments of the actual return block]
\label{thm:exponential-return}
Let
$Q_R=r_R^*+\tau_R^*+|\kappa_R^*|$ on the genuine first-return section.
There are constants $c_1,a_1,C_1>0$, independent of $R\in I$, such that
\begin{equation}\label{eq:exponential-return}
 \nu_R^*(r_R^*>k)\le C_1e^{-c_1k}\quad(k\ge0),\qquad
 \int e^{a_1 Q_R}\dd\nu_R^*\le C_1.
\end{equation}
For $Q_{n,R}=\sum_{j<n}Q_R\circ(F_R^*)^j$ and every $n\ge1$,
\begin{equation}\label{eq:block-exponential-moment}
 \int\exp(a_1 Q_{n,R}/n)\dd\nu_R^*\le C_1,\qquad
 \nu_R^*(N_{n,R}>L)\le C_1e^{-a_1L/n}.
\end{equation}
These are estimates for the original deterministic mechanical records.
\end{theorem}
\begin{proof}
If $x\in Y_R^*$ and $r_R^*(x)>k$, none of $T_Rx,\ldots,T_R^kx$
belongs to $Y_R^*$. Every term of $\sum_{j=1}^k\chi\circ T_R^j$ is
therefore zero. Positivity and invariance of $\nu$ imply
\[
 \nu_R^*(r_R^*>k)
 \le c_*^{-1}\int_M e^{-z_0\sum_{j=1}^k\chi\circ T_R^j}\dd\nu
 \le c_*^{-1}Ce^{-ck}.
\]
This comparison uses neither the regularity of $\one_{Y_R^*}$ as a
multiplier nor an independent-return construction.

Uniform finite horizon and the uniform bound on a single-flight lattice
step give $Q_R\le C_0 r_R^*$ for a fixed $C_0$. Summation of the
exponential tail of the positive integer $r_R^*$ gives a uniform
exponential moment of $Q_R$ for $a_1>0$ sufficiently small. Enlarge
$C_1$ to include the first bound. Finally H\"older's inequality with
$n$ equal exponents and invariance of $F_R^*$ give
\[
 \int\prod_{j<n}\exp\bigl(a_1Q_R\circ(F_R^*)^j/n\bigr)\dd\nu_R^*
 \le\prod_{j<n}\left(\int e^{a_1Q_R\circ(F_R^*)^j}\dd\nu_R^*\right)^{1/n}
 \le C_1.
\]
Since $N_{n,R}\le Q_{n,R}$, exponential Markov inequality proves the
last assertion. No independence is used in this step either.
\end{proof}

\begin{corollary}[All finite moments under parameter variation]
\label{cor:all-moment-continuity}
For fixed $n$ and $1\le p<\infty$, the map
$R\mapsto\widetilde J_{n,R}$ is continuous into $L^p(\nu;\R^4)$.
Every polynomial moment of a fixed finite list of return records is
finite and continuous in $R$. In particular every finite-lag covariance
of the induced record is continuous.
\end{corollary}
\begin{proof}
Almost-everywhere convergence is Lemma~\ref{lem:finite-return-stability}.
Equation~\eqref{eq:block-exponential-moment}, the fixed section mass and
$|J_{n,R}|\le Q_{n,R}$ make all fixed powers uniformly integrable.
Truncation followed by dominated convergence proves $L^p$ convergence
along every convergent parameter sequence. A finite list of records can
be expressed as differences of cumulative records. H\"older's inequality
and the same exponential moment bound give uniform integrability of each
polynomial in that list. Truncation proves its continuity. This proves
finite-lag, not yet infinite Green--Kubo, covariance continuity.
\end{proof}

\section{Exponential control of cumulative collision counts}
\label{sec:cumulative-returns}
The avoidance estimate controls more than a single return. It also
controls the time of the $n$th visit to the section, without separating
the successive return blocks. This observation improves the truncation
budget in \eqref{eq:physical-count-truncation} and supplies exponential
moments on a fixed complex neighborhood for the entire physical record.

Retain $z_0,c,C$ from Proposition~\ref{prop:soft-killing}. Write
$N_{n,R}=\sum_{j<n}r_R^*\circ(F_R^*)^j$, and enlarge a fixed constant
$D\ge1$, if necessary, so that
\begin{equation}\label{eq:cumulative-record-domination}
 |J_{n,R}|_1\le D N_{n,R},\qquad Q_{n,R}\le D N_{n,R}.
\end{equation}
Such a constant exists because each physical flight has uniformly
bounded length and lattice displacement. All probabilities and
expectations in this section are with respect to $\nu_R^*$.

\begin{theorem}[The time of the $n$th actual return]
\label{thm:cumulative-return-tail}
There are constants $A\ge1$, $a,c>0$, independent of $n$, $L$ and $R$,
such that
\begin{equation}\label{eq:cumulative-return-tail}
 \nu_R^*\{N_{n,R}>L\}\le A\exp(an-cL)
       \qquad(n\ge1,\ L\ge0).
\end{equation}
One may take $a=z_0$ and the exponent $c$ of the soft-killing estimate.
Put $v=a/c$. For $0<s<c$ there is a constant $B_s<\infty$, uniform in
$n,R$, for which
\begin{equation}\label{eq:cumulative-exponential-moment}
 \int e^{s(N_{n,R}-vn)_+}\dd\nu_R^*\le B_s,
 \qquad \int e^{sN_{n,R}}\dd\nu_R^*\le B_s e^{svn}.
\end{equation}
Consequently, for every $0<\lambda<c/D$,
\begin{equation}\label{eq:fixed-strip-record-moment}
 \int e^{\lambda Q_{n,R}}\dd\nu_R^*
          \le B_{\lambda D}e^{\lambda Dvn}.
\end{equation}
The constants $B_s$ are locally bounded for $0\le s<c$, with $B_0=1$.
The number $v$ is an upper envelope, not the Kac mean $c_*^{-1}$.
\end{theorem}
\begin{proof}
For an integer $L\ge0$ let
\[
 H_{L,R}(x)=\sum_{j=1}^{L}\one_{Y_R^*}(T_R^j x).
\]
Since $N_{n,R}$ is the physical collision time of the $n$th return,
$N_{n,R}>L$ if and only if $H_{L,R}<n$ on $Y_R^*$. The fixed smooth
weight satisfies $0\le\chi\le\one_{Y_R^*}$. Thus, on this event,
$\sum_{j=1}^{L}\chi(T_R^j x)\le n-1$. In particular, pointwise on
$Y_R^*$,
\[
 \one_{\{N_{n,R}>L\}}
 \le e^{z_0(n-1)}
       \exp\left(-z_0\sum_{j=1}^{L}\chi(T_R^j x)\right).
\]
Integrate on $Y_R^*$, enlarge the domain to $M$, and use invariance to
shift the sum by one collision. Equation~\eqref{eq:soft-killing} gives
\[
 \nu_R^*\{N_{n,R}>L\}\le c_*^{-1}C e^{z_0(n-1)-cL}.
\]
For real $L$ replace it by $\lfloor L\rfloor$ and enlarge the constant
by $e^c$. This proves \eqref{eq:cumulative-return-tail}. It uses a single
orbit's visit count, not a product of probabilities of return blocks.

It follows that, for $u\ge0$,
$\nu_R^*\{N_{n,R}>vn+u\}\le A e^{-cu}$. For a nonnegative random
variable $X$, Tonelli's theorem gives
\[
 \int e^{sX}\dd\nu_R^*
       =1+s\int_0^\infty e^{su}\nu_R^*\{X>u\}\dd u.
\]
Apply this with $X=(N_{n,R}-vn)_+$ to obtain
$B_s=1+As/(c-s)$. Since $N_{n,R}\le vn+X$, the second bound in
\eqref{eq:cumulative-exponential-moment} follows. Finally apply
\eqref{eq:cumulative-record-domination} with $s=\lambda D$.
\end{proof}

We now apply the tail to the actual record measure. For a measurable
insertion $w_{n,R}$ define the finite signed or complex measure
\[
 \mu_{n,R}^{w}=(J_{n,R})_*(w_{n,R}\nu_R^*),\qquad
 \mu_{n,R}^{w,L}=(J_{n,R})_*
       (w_{n,R}\one_{\{N_{n,R}\le L\}}\nu_R^*).
\]
The transform uses the convention $\widehat\mu(\omega)=
\int e^{i\omega\cdot j}\dd\mu(j)$, with
$\omega=(u_1,u_2,s,b)\in\T^3\times\R$.

\begin{proposition}[Weighted finite-band truncation]
\label{prop:linear-count-budget}
Fix a nonnegative integer $d$ and suppose
$|w_{n,R}|\le M(1+Q_{n,R})^d$. For every multi-index $\gamma$ there is a
constant $C_{d,\gamma}$, independent of $n,R,L,w$, such that
\begin{align}\label{eq:weighted-count-truncation}
 \|\mu_{n,R}^{w}-\mu_{n,R}^{w,L}\|_{\TV}
  &\le C_{d,0}M(1+L)^d e^{an-cL},\\
 \sup_{\omega\in\T^3\times\R}
 \left|\partial_\omega^\gamma
   (\widehat\mu_{n,R}^{w}-\widehat\mu_{n,R}^{w,L})(\omega)\right|
  &\le C_{d,\gamma}M(1+L)^{d+|\gamma|}e^{an-cL}.\nonumber
\end{align}
For a bounded insertion ($d=0$), a frequency region of finite volume
$V>0$, and $\varepsilon>0$, the sufficient integer cutoff for integrated
transform error at most $\varepsilon$ is
\begin{equation}\label{eq:linear-count-budget}
 L\ge\left\lceil\frac{a}{c}n+
       \frac1c\max\left\{0,\log\frac{C_{0,0}MV}{\varepsilon}\right\}
        \right\rceil.
\end{equation}
If $V_n\le V_0e^{\kappa n}$ and the required error is $e^{-\delta n}$,
then a cutoff linear in $n$ suffices. This remains true for each fixed
$d,\gamma$, with any fixed positive slack in the exponential budget.
\end{proposition}
\begin{proof}
For an integer $m\ge0$, integration of the tail gives
\begin{align*}
 &\int (1+N_{n,R})^m\one_{\{N_{n,R}>L\}}\dd\nu_R^*\\
 &\quad=(1+L)^m\nu_R^*\{N_{n,R}>L\}
    +m\int_L^\infty(1+t)^{m-1}\nu_R^*\{N_{n,R}>t\}\dd t\\
 &\quad\le C_m(1+L)^m e^{an-cL}.
\end{align*}
The integral term is zero when $m=0$. For $m\ge1$, the last inequality
follows by writing $t=L+u$ and using
$1+L+u\le(1+L)(1+u)$. Pushforward does not increase total variation.
For the derivative estimate differentiate under the integral: its
absolute integrand is at most
$M(1+Q_{n,R})^d|J_{n,R}|_1^{|\gamma|}
\one_{\{N_{n,R}>L\}}$. Equation~\eqref{eq:cumulative-record-domination}
and the preceding tail integral prove both estimates.

Multiplication by $V$ bounds the integral of the transform error. Solving
$C_{0,0}MV e^{an-cL}\le\varepsilon$ gives
\eqref{eq:linear-count-budget}. More generally choose a slope
$\ell>(a+\kappa+\delta)/c$. Then, for $L_n=\lceil\ell n\rceil$,
\[
 V_n(1+L_n)^{d+|\gamma|}e^{an-cL_n}
       \le C(1+n)^{d+|\gamma|}e^{-(\delta+\epsilon_0)n}
\]
for some $\epsilon_0>0$. A fixed additive increase of the cutoffs absorbs
the remaining constant and polynomial, proving the last assertion.
\end{proof}

\begin{remark}
This improves the sufficient cutoff of order
$n\log(V/\varepsilon)$ obtained from
\eqref{eq:physical-count-truncation} to order
$n+\log(V/\varepsilon)$. On an exponentially growing frequency band the
former budget is quadratic in $n$, whereas the new one is linear.
Neither estimate bounds the number of critical words, the variation of
an inverse coarea Jacobian, or a pointwise raw density. In particular the
all-branch derivative sum in Lemma~\ref{lem:raw-residual-sum} remains a
separate requirement. The truncation removes only actual high-count
records in an error estimate; the limiting datum is still the full
unmodified physical record.
\end{remark}

\section{A common realization and the physical renewal identity}
\label{sec:common-renewal}
The preceding moment estimate gives a common realization of the
unbounded record twists on a scale of Lebesgue spaces. It is useful to
construct this realization before seeking an anisotropic spectral
estimate: the collision and induced operators, their damping variables,
and their chronological composition can then be identified without
regularity assumptions on an indicator multiplier.

\subsection{First-return operators on the common collision space}
Use $M=\T\times(-1,1)$ and $\dd\nu=(4\pi)^{-1}\dd\alpha\dd p$.
Let $P_R$ be multiplication by $\one_{Y_R^*}$ and let $Q_R^\perp=I-P_R$.
The superscript distinguishes this projection from the scalar block
size $Q_R$ used earlier. On $L^1(M,\nu)$, define
\[
 (\mathcal A_{R,\omega}h)(y)
   =e^{i\omega\cdot f_R(T_R^{-1}y)}h(T_R^{-1}y),\qquad
 f_R=(\kappa_{{\rm coll},1},\kappa_{{\rm coll},2},1,\tau_R),
\]
for real $\omega=(u_1,u_2,s,b)$. Collision singularities are discarded
as null sets. These are isometries. The operators $P_R,Q_R^\perp$ are
contractions on every $L^p$, although no such assertion about their
action on anisotropic distribution spaces is needed or made here.

\begin{theorem}[Exact damped renewal realization]
\label{thm:common-renewal}
For $m\ge1$ set
\begin{equation}\label{eq:first-return-operator}
 \mathcal R_{m,R,\omega}
 =P_R\mathcal A_{R,\omega}
       (Q_R^\perp\mathcal A_{R,\omega})^{m-1}P_R,
 \qquad
 \mathcal R_R(z,\omega)=\sum_{m\ge1}z^m\mathcal R_{m,R,\omega}.
\end{equation}
For $|z|<1$ the series converges in operator norm on $L^1(M,\nu)$ and
\begin{equation}\label{eq:induced-damped-bound}
 \|\mathcal R_R(z,\omega)h\|_1\le |z|\|P_Rh\|_1.
\end{equation}
On the closed subspace $P_RL^1$, with identity $P_R$, one has
\begin{equation}\label{eq:physical-renewal-identity}
 P_R(I-z\mathcal A_{R,\omega})^{-1}P_R
       =(I-\mathcal R_R(z,\omega))^{-1}.
\end{equation}
Moreover,
\begin{equation}\label{eq:schur-return-operator}
 \begin{split}
 \mathcal R_R(z,\omega)={}&zP_R\mathcal A P_R\\
 &+z^2P_R\mathcal A Q_R^\perp
 (I-zQ_R^\perp\mathcal A Q_R^\perp)^{-1}
             Q_R^\perp\mathcal A P_R,
 \end{split}
\end{equation}
where $\mathcal A=\mathcal A_{R,\omega}$ and the middle inverse acts
on $Q_R^\perp L^1$.

For $|z|\le1$, the strongly defined boundary series represents the
actual first-return twist: for $y\in Y_R^*$ and $x=(F_R^*)^{-1}y$,
\begin{equation}\label{eq:actual-induced-realization}
 (\mathcal R_R(z,\omega)h)(y)
       =z^{r_R^*(x)}e^{i\omega\cdot G_R(x)}h(x).
\end{equation}
It is zero off $Y_R^*$. In particular, for every $n\ge1$,
\begin{equation}\label{eq:actual-record-pairing}
 c_*^{-1}\int_M\mathcal R_R(z,\omega)^n\one_{Y_R^*}\dd\nu
       =\int_{Y_R^*}z^{N_{n,R}}e^{i\omega\cdot J_{n,R}}\dd\nu_R^*.
\end{equation}
\end{theorem}
\begin{proof}
The collision map is invertible and preserves $\nu$. Poincar\'e
recurrence in both time directions makes $F_R^*$ invertible and
measure-preserving on $Y_R^*$, modulo null sets. In the product
\eqref{eq:first-return-operator}, the initial and final states lie in
$Y_R^*$ and the $m-1$ intervening states lie outside it. Thus its initial
domain is exactly $\{r_R^*=m\}$. Its weight is the exponential of the
sum of the $m$ actual collision records, hence $e^{i\omega\cdot G_R}$.
These initial domains partition $Y_R^*$ up to a null set. Their images
under the corresponding return branches also partition $Y_R^*$, by
invertibility of the induced map. Consequently
\[
 \sum_{m\ge1}|z|^m\|\mathcal R_{m,R,\omega}h\|_1
       =\int_{Y_R^*}|z|^{r_R^*}|h|\dd\nu
       \le |z|\|P_Rh\|_1 \quad (|z|\le1).
\]
Each summand has operator norm at most one, so for $|z|<1$ the
operator-norm series converges geometrically. At the boundary the last
display proves strong convergence for each $h$ and gives
\eqref{eq:actual-induced-realization}. It does not assert operator-norm
convergence of the boundary series.

Let $\mathcal U_0=P_R$ and $\mathcal U_m=P_R\mathcal A^mP_R$ for
$m\ge1$. Partition a trajectory with endpoints in $Y_R^*$ according to
its first return time. Chronological composition, with the first segment
on the right, gives
\[
 \mathcal U_m=\sum_{r=1}^m\mathcal U_{m-r}\mathcal R_{r,R,\omega}.
\]
All generating series converge in operator norm for $|z|<1$. Summing
the identity gives $\mathcal U(z)=P_R+\mathcal U(z)\mathcal R_R(z,\omega)$.
The inverse of $I-\mathcal R_R$ exists by
\eqref{eq:induced-damped-bound}, and
$\mathcal U(z)=P_R(I-z\mathcal A)^{-1}P_R$, proving
\eqref{eq:physical-renewal-identity}. Separating $m=1$ from the series
and summing the intervening outside-section iterates proves
\eqref{eq:schur-return-operator}. All these inverses are used only in
$|z|<1$.

Finally iterate \eqref{eq:actual-induced-realization} and change
variables by $(F_R^*)^n$. Its collision counts and record sums are
$N_{n,R}$ and $J_{n,R}$, respectively. Division by $c_*$ yields
\eqref{eq:actual-record-pairing}. No independence of successive
excursions enters the identity.
\end{proof}

\subsection{A fixed complex neighborhood for the unbounded twists}
To cross the damping boundary at $z=1$, use $z=e^w$. Define the operator
directly by the actual $n$th return, rather than by an unproved
continuation of a resolvent. For $y\in Y_R^*$ let $x=(F_R^*)^{-n}y$ and set
\begin{equation}\label{eq:block-scale-operator}
 (\mathcal T_{n,R}(w,\omega)h)(y)
       =e^{wN_{n,R}(x)+i\omega\cdot J_{n,R}(x)}h(x),
\end{equation}
with zero extension to $M\setminus Y_R^*$. Here $w\in\mathbb C$ and
$\omega\in\mathbb C^4$; the first three real frequencies are interpreted
modulo $2\pi$. The symbols $J_{n,R}$ and $G_R$ denote four real
coordinates, with the displacement counted twice.

\begin{theorem}[Holomorphic realization on a common Lebesgue scale]
\label{thm:common-scale-holomorphy}
Fix $1\le q<p\le\infty$ and put $t=(q^{-1}-p^{-1})^{-1}$. Let
\[
 \sigma(w,\omega)=|\operatorname{Re}w|+D|\operatorname{Im}\omega|_\infty,
 \qquad \Omega_{p,q}=\{(w,\omega):t\sigma(w,\omega)<c\}.
\]
For every $n\ge1$, \eqref{eq:block-scale-operator} defines a
holomorphic map from $\Omega_{p,q}$ to
$\mathcal L(L^p(M,\nu),L^q(M,\nu))$. Its norm satisfies
\begin{equation}\label{eq:scale-operator-bound}
 \|\mathcal T_{n,R}(w,\omega)\|_{p\to q}
       \le c_*^{1/t}B_{t\sigma}^{1/t}e^{\sigma vn}.
\end{equation}
On every compact subset of $\Omega_{p,q}$, all fixed-order derivatives in $(w,\omega)$ have bounds $C_k e^{C_k n}$ uniform in $R$. At every fixed
$(w,\omega)\in\Omega_{p,q}$ and every $h\in L^p(M,\nu)$,
\begin{equation}\label{eq:scale-parameter-continuity}
 R\longmapsto\mathcal T_{n,R}(w,\omega)h
       \quad\hbox{is continuous in }L^q(M,\nu).
\end{equation}
For $\operatorname{Re}w<0$ and real $\omega$, this operator agrees
with $\mathcal R_R(e^w,\omega)^n$ wherever the latter is considered
as an $L^1$ operator. The resulting inserted record pairings extend
holomorphically across $w=0$ on this fixed neighborhood.
\end{theorem}
\begin{proof}
Change variables $y=(F_R^*)^n x$ in the $q$th power of the norm. The
absolute weight is at most $e^{\sigma N_{n,R}(x)}$, by
\eqref{eq:cumulative-record-domination}. H\"older's inequality with
$1/q=1/p+1/t$ and Theorem~\ref{thm:cumulative-return-tail} give
\[
 \|\mathcal T_{n,R}(w,\omega)h\|_q
 \le\left(\int_{Y_R^*}e^{t\sigma N_{n,R}}\dd\nu\right)^{1/t}
       \|h\|_p
 \le c_*^{1/t}B_{t\sigma}^{1/t}e^{\sigma vn}\|h\|_p.
\]
For $\sigma=0$ use $B_0=1$. This proves the operator bound, including
$p=\infty$.

On a compact subset of the stated domain choose a slightly larger
compact neighborhood for which $t\sigma'<c$. A derivative of order
$k$ in $(w,\omega)$ inserts a monomial in $N_{n,R}$ and the four
coordinates of $J_{n,R}$, bounded by $C_kN_{n,R}^k$. Every such
polynomial is absorbed by $e^{(\sigma'-\sigma)N_{n,R}}$. H\"older's
inequality and \eqref{eq:cumulative-exponential-moment} therefore bound
the derivative integrands and the Taylor remainders in operator norm.
The exponential power series on a sufficiently small polydisc converges
in that norm. This proves joint holomorphy, not just weak holomorphy of
individual pairings. The same estimates on the larger neighborhood,
or Cauchy's inequalities there, give the claimed derivative bounds.

We give the parameter-continuity argument because the underlying return
domains move. Fix $R_0$. Remove $\partial Y_{R_0}^*$, the forward and
backward collision singularities, and all finite collision iterates of
$\partial Y_{R_0}^*$. Each is a countable union of null curves, up to
singular endpoints. Backward recurrence holds almost everywhere. At a
remaining point of $Y_{R_0}^*$ the preceding $n$ returns use finitely many
nonsingular physical collisions; none is on a section boundary. The
implicit collision-root argument of
Theorem~\ref{thm:return-stability}, applied backwards, preserves this
finite itinerary for $R$ sufficiently near $R_0$. The backward starting
point and all flight lengths vary continuously; the lattice labels and
collision counts are fixed. At a remaining point outside $Y_{R_0}^*$
the section indicator stays zero. Thus
\eqref{eq:block-scale-operator} converges pointwise almost everywhere
when $h$ is bounded and continuous.

Choose $q'>q$ sufficiently close to $q$ so that $q'<p$ and
$t'\sigma<c$, where $t'=(1/q'-1/p)^{-1}$. The preceding bound in
$L^{q'}$ is uniform in nearby $R$. It gives uniform integrability of the
$q$th powers, so pointwise convergence implies convergence in $L^q$.
For $p<\infty$, approximate an arbitrary $h$ in $L^p$ by bounded smooth
functions and use the uniform $p\to q$ bound. For $p=\infty$, choose a
finite $p_0>q$ so large that $(1/q-1/p_0)^{-1}\sigma<c$. Approximation
in $L^{p_0}$ and its uniform $p_0\to q$ bound give the same conclusion
for $h\in L^\infty$. This avoids assuming norm density of smooth
functions in $L^\infty$.

In the damped region the two operators have the same pointwise formula,
by Theorem~\ref{thm:common-renewal}. For example, for $h\in L^p$ and
$b\in L^{q^\dagger}$, where $q^\dagger$ is the conjugate exponent of $q$,
\[
 c_*^{-1}\int b\,\mathcal T_{n,R}(w,\omega)h\dd\nu
 =\int_{Y_R^*}h(x)b((F_R^*)^nx)
       e^{wN_{n,R}(x)+i\omega\cdot J_{n,R}(x)}\dd\nu_R^*.
\]
This continuous pairing is holomorphic on $\Omega_{p,q}$. It agrees
with the damped renewal pairing there and hence is the asserted
extension across $w=0$.
\end{proof}

\begin{remark}[Relation to the spectral and raw-density estimates]
The operator identities concern the genuine four-coordinate record,
and all their spaces live on the same collision probability space.
They specify a common realization with fixed exponential integrability
and strong parameter continuity. They do not assert operator-norm
continuity in $R$, quasi-compactness on $L^p$, or a holomorphic family of
endomorphisms of a single anisotropic space. At real frequencies the
induced $L^1$ operator is an isometry on its section, so its construction
alone gives no spectral gap. Likewise, maps $L^p\to L^q$ with $q<p$
cannot simply be iterated on one space outside the damped region; the
$n$-block operator above was defined directly.

The next analytical step is to establish an anisotropic realization
compatible with \eqref{eq:actual-record-pairing} and
\eqref{eq:physical-renewal-identity}, and to prove its phase
reconstruction, covariance and raw-residual bounds. The criteria in
Theorem~\ref{thm:phase-resolvent-criterion},
Proposition~\ref{prop:covariance-closure}, and
Lemma~\ref{lem:raw-residual-sum} continue to state these requirements
separately. In particular the common Lebesgue realization is not used
as a substitute for any of those estimates in the raw local limit.
\end{remark}

\section{Common charts for moving return domains}
\label{sec:common-charts}
We give a local change-of-variables estimate and then apply it to the
actual return domains. This makes explicit the domination used in
Theorem~\ref{thm:return-stability}. Critical neighborhoods are discarded
in the initial collision space, not by assuming a uniform density bound
near their images.

\begin{lemma}[A common submersion chart]
\label{lem:common-submersion-chart}
Let $U=U_1\times U_2\Subset\R^2$ be a bounded open rectangle and let
$f(R,u,v)$ be $C^2$ on a neighborhood of $J\times\overline U$, where
$J$ is a compact parameter interval. Suppose
$\partial_v f\ge\eta>0$ there. For $a\in C_c^1(U)$, let $g_R$ be the
density of the pushforward of $a(u,v)\dd u\dd v$ by $f_R$. Then
\begin{equation}\label{eq:chart-Lipschitz}
 \|g_R-g_S\|_{L^1(\R)}\le C_a|R-S|\qquad(R,S\in J).
\end{equation}
The constant depends only on $U$, a common bounded range of $f$, $\eta$,
the $C^2$ bound for $f$, and the $C^1$ bound and support of $a$.
The same assertion holds if $\partial_v f\le-\eta$.
\end{lemma}
\begin{proof}
For fixed $R,u$ the map $v\mapsto f(R,u,v)$ is one-to-one. Write its
inverse as $v_R(u,t)$. The pushforward under $(u,v)\mapsto(u,f_R(u,v))$
has density
\[
 A_R(u,t)=\frac{a(u,v_R(u,t))}{\partial_v f(R,u,v_R(u,t))}
\]
on its image, extended by zero elsewhere. All images lie in one bounded
$(u,t)$ box. Since $a$ vanishes on a fixed neighborhood of $\partial U$,
this zero extension introduces no boundary jump as $R$ varies. Wherever
$a$ or its derivatives are nonzero, the inverse is defined on a common
parameter neighborhood. If $M$ bounds the relevant derivatives of $f$,
implicit differentiation gives
\[
 |\partial_R v_R|\le M/\eta,\qquad
 |\partial_R A_R|
 \le \|a_v\|_\infty M/\eta^2
       +\|a\|_\infty(M/\eta^2+M^2/\eta^3).
\]
At points outside the inverse chart both sides are interpreted by the
zero extension, which vanishes on an open neighborhood of its moving
edge. The same bound therefore gives an $L^1$ majorant on the fixed box.
Integrating in $R$ proves
$\|A_R-A_S\|_1\le C_a|R-S|$. Since
$g_R(t)=\int A_R(u,t)\dd u$, integration in $u$ proves
\eqref{eq:chart-Lipschitz}. Reversing the $v$ coordinate handles the
negative derivative. The result also applies on a constant lattice
component of a mixed density.
\end{proof}

\begin{proposition}[Finite common patches with an explicit remainder]
\label{prop:finite-common-patches}
Fix $n\ge1$, $R_0\in I$, and $\epsilon>0$. There are finitely many
nonnegative smooth cutoffs $\psi_i$ on the common collision space, a
parameter neighborhood $J$ of $R_0$, and $C_{\epsilon,n,R_0}<\infty$
such that $\sum_i\psi_i\le1$, every support lies in one common regular
first-$n$-return patch for all $R\in J$, and
\begin{equation}\label{eq:common-patch-mass}
 c_*^{-1}\int\sum_i\psi_i\dd\nu>1-\epsilon.
\end{equation}
The lattice labels and all intermediate return decisions are fixed on
each patch. If $p^i_{n,R}$ is the pushforward density of
$c_*^{-1}\psi_i\nu$, then
\begin{align}\label{eq:common-patch-bound}
 \sum_i\|p^i_{n,R}-p^i_{n,S}\|_1
       &\le C_{\epsilon,n,R_0}|R-S|,\\
 \|p_{n,R}-p_{n,S}\|_1
       &\le2\epsilon+C_{\epsilon,n,R_0}|R-S|
               \qquad(R,S\in J).\nonumber
\end{align}
Before choosing these patches one may discard $N_{n,R_0}>L$ at cost at
most $C_1e^{-a_1L/n}$, with the constants of
Theorem~\ref{thm:exponential-return}.
\end{proposition}
\begin{proof}
The exponential moment bound proves the last assertion. Choose $L$ so
that this cost is less than $\epsilon/3$. On the remainder remove the
initial boundary of the section, all grazing and competing-root
singularities encountered before collision $L$, and the preimages of the
moving return boundary at $R_0$ through those collisions. These sets have
zero collision measure. On every remaining regular word also remove the
critical points of its cumulative flight length. Proposition
\ref{prop:general-edge} shows that there is at most one such point on a
regular optical word; the submersion complement has full measure.
There are only finitely many possible physical lattice words of length
at most $L$, since each flight label belongs to a fixed finite set.
No estimate for the number of unbounded-length words is used here.

Inner regularity of the collision measure and a countable exhaustion of
the regular submersion patches give a compact subset of mass greater
than $1-\epsilon$ in the normalized section. At every point of this
compact subset, choose a small coordinate rectangle on which one
coordinate derivative of the cumulative time is bounded away from zero.
The finite-itinerary lemma preserves the entire physical word and every
return decision on a slightly larger rectangle for $R$ close to $R_0$.
Select finitely many such rectangles and a subordinate family of smooth
cutoffs with sum at most one, equal to one on the compact subset.
Their supports have positive distance from all the excluded sets at
$R_0$. Finiteness gives a common parameter neighborhood $J$ on which
these distances, return decisions and derivative bounds persist. The
time functions and their first two derivatives are uniformly bounded
on these finitely many compact rectangles. Lemma
\ref{lem:common-submersion-chart} applies with
$a=\psi_i/(4\pi c_*)$ in $(\alpha,p)$ coordinates. Summing its bounds
proves the first line of \eqref{eq:common-patch-bound}.

For every $R\in J$ the remainder of the original probability is a
positive measure. Its mass is exactly
$1-c_*^{-1}\int\sum_i\psi_i\dd\nu<\epsilon$, because the section
measure is the fixed number $c_*$. The raw densities exist by the
submersion/coarea argument of Theorem~\ref{thm:return-stability}.
The two remainder densities have total $L^1$ norm less than
$2\epsilon$. This gives the second line. Nothing has been assumed about
the size of a density at a critical value or a sharp image boundary.
\end{proof}

The constants in \eqref{eq:common-patch-bound} expose the order of
choices: physical-count truncation, deletion of small initial-state
neighborhoods, compact common charts, and finally parameter tolerance.
They are not asserted to remain bounded as $n$ or $\epsilon^{-1}$ grows.
A quantitative global density modulus requires quantitative control of
those additional choices. Bounded measurable insertions are handled by
$L^1$ approximation on the same patches, as in
Proposition~\ref{prop:weighted-return-stability}; no artificial observation
coordinate is introduced.

\section{A uniform functional law for the mechanical clock}
\label{sec:uniform-fluid-clock}
Finite-record continuity also has a long-time consequence that does not
require a mixing rate. We prove a parameter-uniform functional law at the
first-order scale, including the maximum unfinished return over an entire
observation interval. The \(\sqrt t\)-scale process theorem in
Proposition~\ref{prop:clock} remains a separate statement.

Set
\[
 G_R=(\kappa_R^*,r_R^*,\tau_R^*),\qquad
 \bar G_R=(0,0,c_*^{-1},\bar\tau_R/c_*),\qquad J_{0,R}=0.
\]
All probabilities on the return section use $\nu_R^*$, not an
independent sequence of records. Write $\bar r=c_*^{-1}$ and
$\bar\tau_R^*=\bar\tau_R/c_*$. Both $\bar\tau_R^*$ and its reciprocal
are bounded on $I$.

\begin{lemma}[A compact subadditive argument]
\label{lem:compact-subadditive}
Let $a_n:I\to[0,\infty)$ be continuous, $a_0=0$, and
$a_{n+m}(R)\le a_n(R)+a_m(R)$. Suppose
$\inf_{n\ge1}a_n(R)/n=0$ for every $R$. Then
\[
 \lim_{n\to\infty}\sup_{R\in I}\frac{a_n(R)}n=0.
\]
More precisely, for every $\epsilon>0$ there is a finite integer $M$
such that, with $C=\sup_R a_1(R)$,
\begin{equation}\label{eq:subadditive-uniform-budget}
 \sup_R a_n(R)/n\le\epsilon+CM/n\qquad(n\ge1).
\end{equation}
\end{lemma}
\begin{proof}
For each $R$ choose a block length $m_R$ with
$a_{m_R}(R)/m_R<\epsilon$. Continuity preserves this strict inequality
on a neighborhood of $R$. Take a finite subcover and let $M$ be the
largest of its block lengths. For an arbitrary parameter choose one of
the covering neighborhoods, with block length $m\le M$, and write
$n=qm+r$, $0\le r<m$. Subadditivity and $a_r\le r a_1$ give
$a_n\le q\epsilon m+Cr\le\epsilon n+CM$. This proves the
estimate and then the limit. The neighborhoods need not have a common
block length; replacing them by an unproved common mixing time would
not be justified.
\end{proof}

\begin{theorem}[Uniform first-order return process]
\label{thm:uniform-return-fluid}
The actual induced records satisfy
\begin{equation}\label{eq:uniform-return-mean-law}
 \lim_{n\to\infty}\sup_{R\in I}
       \int\left|J_{n,R}/n-\bar G_R\right|\dd\nu_R^*=0.
\end{equation}
For every fixed $A<\infty$ and $\epsilon>0$,
\begin{equation}\label{eq:uniform-return-functional}
 \lim_{n\to\infty}\sup_{R\in I}\nu_R^*\left\{
 \sup_{0\le u\le A}
 \left|n^{-1}J_{\lfloor nu\rfloor,R}-u\bar G_R\right|>\epsilon
 \right\}=0.
\end{equation}
\end{theorem}
\begin{proof}
The ergodic collision map in Section~\ref{sec:return-stability}
induces an ergodic map on its positive-measure return section. Indeed,
the complete tower over an invariant subset of the induced map is an
invariant collision set. The two towers over that subset and its
complement partition the full collision space modulo null sets;
collision ergodicity forces one of the two to have zero measure.
Thus the integrable vector $G_R$ obeys the $L^1$ ergodic theorem for
each fixed $R$, with the exact mean in
Proposition~\ref{prop:exact-induced-means}.

Put $a_n(R)=\int|J_{n,R}-n\bar G_R|\dd\nu_R^*$.
Invariance of $F_R^*$ and the cocycle identity make $a_n$ subadditive.
It is continuous: on the common collision space it equals
\[
 c_*^{-1}\int_M
 \left|\widetilde J_{n,R}-n\bar G_R\one_{Y_R^*}\right|\dd\nu.
\]
Theorem~\ref{thm:return-stability} gives the first term's $L^1$
continuity, the section indicators are continuous in $L^1$ because
their boundaries move continuously and have zero measure, and the
mean is continuous by its explicit formula. The pointwise ergodic
theorem gives $a_n(R)/n\to0$. Lemma~\ref{lem:compact-subadditive}
proves \eqref{eq:uniform-return-mean-law}.

For the functional conclusion, first test a fixed finite mesh of
$[0,A]$. Equation~\eqref{eq:uniform-return-mean-law} and Markov's
inequality give convergence in probability at all its mesh points,
uniformly in $R$. The count and roof coordinates of $J_n$ are
nondecreasing. Between mesh points their normalized deviations are
therefore bounded by the endpoint deviations plus a uniformly bounded
mean times the mesh width, with an $O(n^{-1})$ integer-part error.
For the displacement coordinate, finite horizon gives a constant
$C_\kappa$ such that for every $l\ge j$,
\[
 |K_{l,R}-K_{j,R}|\le C_\kappa(N_{l,R}-N_{j,R}).
\]
Consequently its oscillation on a mesh interval is controlled by the
count increment on that interval. Its mean is zero. On the event
where all mesh errors are small, every coordinate has a uniformly
small error between mesh points as well. Let $n$ grow, then let the
mesh width and the mesh error tolerance decrease. This proves
\eqref{eq:uniform-return-functional} without an independence or
maximal-martingale assumption.
\end{proof}

\begin{lemma}[Sublinear maximum return blocks]
\label{lem:uniform-maximal-block}
Put $Q_R=r_R^*+\tau_R^*+|\kappa_R^*|$ and
\[
 U(L)=\sup_R\int Q_R\one_{\{Q_R>L\}}\dd\nu_R^*.
\]
Then $U(L)\to0$. For $n\ge1$, $A\ge0$, and $\epsilon>0$,
\begin{equation}\label{eq:uniform-maximal-block}
 \sup_R\nu_R^*\left\{
 \max_{0\le j\le\lfloor An\rfloor}Q_R\circ(F_R^*)^j>\epsilon n
 \right\}\le\frac{A+1}{\epsilon}U(\epsilon n).
\end{equation}
Both \eqref{eq:uniform-return-functional} and the convergence to zero
in \eqref{eq:uniform-maximal-block} remain valid when the initial
base point has the stationary length-biased law
\begin{equation}\label{eq:length-biased-base-v4}
 \dd\rho_R=\frac{\tau_R^*}{\bar\tau_R^*}\dd\nu_R^*.
\end{equation}
\end{lemma}
\begin{proof}
Finite horizon bounds $\tau_R^*$ and $|\kappa_R^*|$ by a uniform
constant times $r_R^*$. The uniform integrability in
Theorem~\ref{thm:return-stability}, and the constant section mass,
therefore give $U(L)\to0$. Invariance of $F_R^*$, the union bound
and truncated Markov inequality give
\[
 \nu_R^*\{\max_{j\le\lfloor An\rfloor}Q_R\circ F^j>\epsilon n\}
 \le (\lfloor An\rfloor+1)\nu_R^*(Q_R>\epsilon n)
 \le\frac{A+1}{\epsilon}\int Q_R\one_{Q_R>\epsilon n}\dd\nu_R^*.
\]
No independence of the events is used.

Let $d_*:=\inf_R\bar\tau_R^*>0$. For any event $B_R$ determined
by the forward base orbit and every $L>0$,
\begin{equation}\label{eq:length-bias-transfer-v4}
 \rho_R(B_R)\le d_*^{-1}\left(
 L\nu_R^*(B_R)+\int\tau_R^*\one_{\{\tau_R^*>L\}}\dd\nu_R^*\right).
\end{equation}
If $\sup_R\nu_R^*(B_R)\to0$, first hold $L$ fixed and pass to that
limit, then let $L\to\infty$. Uniform integrability of the roof
makes the right side tend to zero uniformly. Apply this argument to
the exceptional events in the two displayed convergence statements.
This transfer uses a length-biased law, not the unweighted launch law.
\end{proof}

For the stationary physical flow, let $V_R(t)$ count future visits to
$Y_R^*$, let $C_R(t)$ count future physical collisions, and let
$W_R(t)\in\Z^2$ be the net fundamental-cell displacement since time
zero. Changing the choice of a fixed fundamental cell changes the
continuous-position comparison only by a bounded term.

\begin{theorem}[Uniform stationary physical clock]
\label{thm:uniform-physical-clock}
For every $\epsilon>0$, under the actual stationary billiard laws
$\mu_R$, uniformly for $R\in I$,
\begin{equation}\label{eq:uniform-physical-clock}
 \sup_{0\le u\le1}\left|
 \frac1t\bigl(V_R(tu),C_R(tu),W_R(tu)\bigr)
 -\left(\frac{uc_*}{\bar\tau_R},
        \frac{u}{\bar\tau_R},0\right)\right|
 \longrightarrow0
 \quad\hbox{in probability as }t\to\infty.
\end{equation}
The maximum marked unfinished return encountered during $[0,t]$,
including the initial return containing time zero, is $o_P(t)$
uniformly in $R$.
\end{theorem}
\begin{proof}
Represent the starting state by $(x,s)$ in the induced suspension,
$0\le s<\tau_R^*(x)$. The marginal law of $x$ is exactly $\rho_R$
in \eqref{eq:length-biased-base-v4}. Write $T_{j,R}$ for the roof
coordinate of $J_{j,R}$. Count future events in $(0,t]$. With this
convention the following inverse identity also holds at event times:
\[
 V_R(tu)=\max\{j\ge0:T_{j,R}(x)\le s+tu\}.
\]
Choose a fixed $A$ with $A d_*>2$. The functional return law,
transferred by Lemma~\ref{lem:uniform-maximal-block}, shows that
$T_{\lfloor At\rfloor,R}>s+t$ with probability tending to one
uniformly. Here $s/t\to0$ uniformly because $s\le Q_R(x)$ and
the same lemma applies at $j=0$. Integer parts do not affect these
statements; the return-process theorem may be applied with
$n=\lceil t\rceil$ and a slightly larger fixed horizon.

On this event all inverse indices lie below $At$. Uniformly for
$0\le j\le At$ we have
$|T_{j,R}-j\bar\tau_R^*|/t=o_P(1)$ and
$\max_{j\le At}Q_R\circ F^j/t=o_P(1)$. The inequalities
\[
 T_{V_R(tu),R}\le s+tu<T_{V_R(tu)+1,R}
\]
therefore imply
\[
 \sup_{u\le1}\left|
 \bar\tau_R^* V_R(tu)/t-u\right|=o_P(1)
\]
uniformly in $R$. Division by $\bar\tau_R^*\ge d_*$ yields the
first component in \eqref{eq:uniform-physical-clock}.

The completed-return count $N_{V_R(tu),R}$ differs from the actual
physical count by the initial prefix and the final incomplete return.
Their sum is at most a fixed constant plus twice the maximal $Q_R$
over these return indices. The displacement comparison has the same
bound, up to a uniform multiplicative constant, since each physical
lattice step is bounded. These bounds hold simultaneously for all
$u\in[0,1]$. The functional return law at the random inverse indices
now gives
\[
 C_R(tu)/t=\bar r\,V_R(tu)/t+o_P(1),\qquad
 W_R(tu)/t=o_P(1)
\]
with uniform suprema. Since $\bar r/\bar\tau_R^*=1/\bar\tau_R$,
this proves the remaining components. Every unfinished piece in
$[0,t]$ belongs to one of the same return blocks and is bounded by
its marked size, proving the final assertion.
\end{proof}

\begin{corollary}[Geometric error in the physical collision rate]
\label{cor:physical-rate-error}
Let $b_*:=\min_{R\in I}\bar\tau_R$ and put
$L_*=15/(4b_*^2)$. There is a function $\omega_t(\epsilon)\to0$
for every fixed $\epsilon>0$ such that, for any $I$-valued radius
estimate $\widehat R$ on a probability space carrying the stationary
billiard at the true radius $R$, and any $\delta>0$,
\begin{align}\label{eq:physical-rate-error}
 \Pr_R\left\{\sup_{0\le u\le1}
 \left|C_R(tu)/t-u/\bar\tau_{\widehat R}\right|
                       >\epsilon+L_*\delta\right\}
 \le\omega_t(\epsilon)+\Pr_R\{|\widehat R-R|>\delta\}.
\end{align}
No independence of $\widehat R$ and the observed orbit is required.
For the visit rate, replace the deterministic error constant by $c_*L_*$.
\end{corollary}
\begin{proof}
The explicit mean satisfies $|\bar\tau_R'|<15/4$, so
$|(1/\bar\tau_R)'|\le L_*$. On the event
$|\widehat R-R|\le\delta$ the two proposed rates differ by at
most $L_*\delta$, uniformly over $u\in[0,1]$. The triangle
inequality and Theorem~\ref{thm:uniform-physical-clock} prove the
claim with $\omega_t$ equal to the supremum, over $R$, of the true
rate's exceptional probability. Multiplication by $c_*$ gives the
visit-rate statement.
\end{proof}

The supremum statement here is at scale $t$. Uniform integrability
alone does not turn it into an $o_P(\sqrt t)$ bound for a growing
window. Likewise \eqref{eq:subadditive-uniform-budget} and
\eqref{eq:physical-rate-error} do not supply a numerical convergence
rate for $\omega_t$ without quantitative finite-block estimates.
They establish the full first-order clock comparison for this moving
family; the covariance, Gaussian process, raw density and shrinking-
conditioning conclusions retain their stronger hypotheses. The
stationary law throughout is not substituted for an unknown experimental
preparation law.

\section{A logarithmic unfinished-block bound under conditioning}
\label{sec:conditioned-clock}
The stationary law of an induced block is length-biased. We use this law
for the initial block and the exact visit intensity for subsequent blocks.
These two ingredients avoid applying a launch-law union bound to a
stationary initial state.

Write $\overline\tau_R^*=\bar\tau_R/c_*$ and represent the flow as the
suspension over $(Y_R^*,F_R^*,\nu_R^*)$ with roof $\tau_R^*$. Let $t_j$
be its two-sided visit times and $x_j$ their outgoing section states.
The minimum time between consecutive visits is at least $3/50$.

\begin{lemma}[Marked visit intensity]
\label{lem:marked-visit-intensity}
For every nonnegative measurable $f$ on $Y_R^*$ and every $t\ge0$,
\begin{equation}\label{eq:marked-visit-intensity}
 \mathbb E_{\mu_R}\sum_{t_j\in(0,t]} f(x_j)
       =\frac{t}{\overline\tau_R^*}\int f\dd\nu_R^*.
\end{equation}
\end{lemma}
\begin{proof}
First take bounded $f$. Stationarity makes the locally finite measure
$A\mapsto\mathbb E_{\mu_R}\sum_{t_j\in A}f(x_j)$ translation-invariant
on $\R$, hence equal to $i_f$ times Lebesgue measure. Local finiteness
follows from the positive lower roof bound. Choose $0<\delta<3/50$.
At a stationary time, being in the first $\delta$ units of an induced
roof is the same event as a visit in the preceding interval of length
$\delta$, apart from null endpoints. There is at most one such visit,
and its mark is the current base state. Integration in the suspension
therefore gives
\[
 \delta i_f=\frac1{\overline\tau_R^*}
       \int_{Y_R^*}\int_0^\delta f(x)\dd s\dd\nu_R^*(x)
       =\frac{\delta}{\overline\tau_R^*}\int f\dd\nu_R^*.
\]
This determines $i_f$ and proves the formula for bounded $f$.
Monotone convergence extends it to arbitrary nonnegative marks, including
the case of an infinite integral. No renewal independence is involved.
\end{proof}

Let $\mathcal M_R(t)$ be the maximum of $Q_R$ over all complete induced
blocks whose time intervals meet $[0,t]$, including the block containing
time zero. Endpoint conventions at a visit are immaterial for $\mu_R$.
The marked size of any unfinished piece is bounded by a fixed multiple
of $1+\mathcal M_R(t)$.

\begin{theorem}[Uniform exponential maximum and conditioned residual]
\label{thm:conditioned-maximal-block}
There are constants $a_2,C_2>0$ such that, for every $R\in I$, $t\ge0$
and $q\ge0$,
\begin{equation}\label{eq:exponential-maximal-block}
 \mu_R\{\mathcal M_R(t)>q\}\le C_2(1+t)e^{-a_2q}.
\end{equation}
Consequently, for any measurable event $B_{R,t}$ of probability at least
$d_t>0$,
\begin{equation}\label{eq:conditioned-maximal-block}
 \mu_R\{\mathcal M_R(t)>q\mid B_{R,t}\}
       \le\min\{1,C_2(1+t)e^{-a_2q}/d_t\}.
\end{equation}
In particular, if $d_t\ge d(1+t)^{-A}$ with $d>0$ and $A\ge0$, then
$\mathcal M_R(t)=O_P(\log(2+t))$ under these conditional laws, uniformly
in $R$ and in the events. The maximum marked unfinished return during
$[0,t]$ is then $o_P(\sqrt t)$ under the same laws. More generally the
last assertion holds if $\log((1+t)/d_t)=o(\sqrt t)$.
\end{theorem}
\begin{proof}
The initial containing block has law
$\tau_R^*\dd\nu_R^*/\overline\tau_R^*$. If the maximum exceeds $q$,
either that block has $Q_R>q$ or a visit in $(0,t]$ starts a block with
$Q_R>q$. The union bound and Lemma~\ref{lem:marked-visit-intensity} give
\begin{equation}\label{eq:exact-maximal-tail-budget}
 \mu_R\{\mathcal M_R(t)>q\}
 \le\frac1{\overline\tau_R^*}
       \left(\int\tau_R^*\one_{\{Q_R>q\}}\dd\nu_R^*
                       +t\,\nu_R^*(Q_R>q)\right).
\end{equation}
The denominator is bounded below on $I$. Since $\tau_R^*\le Q_R$,
the exponential moment in \eqref{eq:exponential-return} bounds the first
integral by $C e^{-a_1q/2}$ and the probability by $C e^{-a_1q}$.
For example, $Q\le (2/a_1)e^{a_1Q/2}$ and
$e^{a_1Q/2}\one_{Q>q}\le e^{a_1Q}e^{-a_1q/2}$ suffice for the
first bound. This proves \eqref{eq:exponential-maximal-block}.

Dividing the unconditional exceptional probability by $\mu_R(B_{R,t})$
proves \eqref{eq:conditioned-maximal-block}. With a polynomial lower
bound on $d_t$, setting $q=K\log(2+t)$ and taking
$a_2K>A+1$ makes the right side tend to zero. With the more general
hypothesis, set $q=\epsilon\sqrt t$ for any $\epsilon>0$; the logarithm
of the displayed upper bound tends to minus infinity. The deterministic
bound of each marked piece by $C(1+Q_R)$ proves the assertions for all
unfinished pieces at once. The event may depend on the whole orbit and
on an estimator from that orbit; no independence is assumed.
\end{proof}

For a stationary starting point $(x,s)$ put
$V_R(v)=\max\{j:T_{j,R}(x)\le s+v\}$. Define the two path records
\[
 Z_R(v)=(V_R(v),C_R(v),W_R(v)),\qquad
 Z_R^\circ(v)=(V_R(v),N_{V_R(v),R}(x),K_{V_R(v),R}(x)).
\]
The second record counts full blocks from the beginning of the initial
containing roof, not from the observation time. It is used solely as an
exact intermediate record for comparison.

\begin{corollary}[A conditional path-comparison budget]
\label{cor:conditional-path-budget}
There is $C_3$ such that, pathwise outside the collision singularities,
\begin{equation}\label{eq:deterministic-prefix-bound}
 \sup_{0\le v\le t}|Z_R(v)-Z_R^\circ(v)|
                         \le C_3(1+\mathcal M_R(t)).
\end{equation}
Let $\widehat Z_R$ and $\widehat Z_R^\circ$ be the corresponding paths
on $[0,1]$, divided by $\sqrt t$ and with any identical deterministic
centering. For $|F|\le1$ Lipschitz with constant $L_F$ for the uniform
path metric and $\mu_R(B_{R,t})\ge d_t>0$,
\begin{align}\label{eq:conditional-path-budget}
 \left|\mathbb E_{\mu_R}[F(\widehat Z_R)\mid B_{R,t}]
       -\mathbb E_{\mu_R}[F(\widehat Z_R^\circ)\mid B_{R,t}]\right|
 &\le\frac{C_3L_F(1+q)}{\sqrt t}
        +\frac{2C_2(1+t)e^{-a_2q}}{d_t}.
\end{align}
For $t\ge2$ choose
\begin{equation}\label{eq:conditional-log-budget}
 q=a_2^{-1}\log\left(\frac{\max\{1,2C_2\}(1+t)\sqrt t}{d_t}\right).
\end{equation}
The second term is at most $t^{-1/2}$; for polynomially bounded
conditioning probabilities the total error is
$O((1+L_F)\log t/\sqrt t)$.
\end{corollary}
\begin{proof}
The collision counts differ by the initial prefix and the final
incomplete block. Each is bounded by the full block's collision count.
The same holds for the lattice displacement after multiplying by the
uniform bound on a single-flight lattice step and adding a fixed
fundamental-cell boundary constant. The visit counts agree exactly by
construction. These estimates are simultaneous for all $v\in[0,t]$,
and prove \eqref{eq:deterministic-prefix-bound}. On
$\{\mathcal M_R(t)\le q\}$ the Lipschitz bound gives the first term in
\eqref{eq:conditional-path-budget}. On its complement the difference of
two values of $F$ is at most two. Apply
\eqref{eq:conditioned-maximal-block}. The stated choice of $q$ and
substitution prove the final assertions.
\end{proof}

\begin{remark}[Which conditional interface has been closed]
\label{rem:conditional-interface-scope}
The square-root unfinished-block requirement is now a theorem for the
actual mechanical family, even after any conditioning of polynomially
small positive probability and over the whole observation interval. A
functional central limit theorem for completed records is still a
separate premise in Proposition~\ref{prop:clock}. For a local limit, the
conditioning event on the two sides of \eqref{eq:conditional-path-budget}
is the same event. Replacing an exact lattice observation by the
completed-block observation changes that event and needs a further
boundary estimate. If two events $B,B'$ have probabilities at least $d$,
then for $|f|\le1$ direct subtraction of the ratios gives
\[
 |\mathbb E[f\mid B]-\mathbb E[f\mid B']|
             \le2\mu_R(B\mathbin\triangle B')/d.
\]
Thus an $O(\log t)$ error in an integer record is not discarded when
conditioning on one lattice value. The weighted raw-density estimates
of Proposition~\ref{prop:conditioning} remain the requisite local input.
A zero-probability singleton in the continuous coordinate is not an event
to which \eqref{eq:conditioned-maximal-block} applies.
\end{remark}

\section{The remaining covariance and raw-summation interfaces}
\label{sec:closure-contracts}
We record the exact analytical implications needed in the low-frequency
and outer-frequency parts of the original density theorem. These
criteria distinguish proved finite-record facts from the two infinite
sums that still require estimates.

\begin{proposition}[A covariance closure criterion]
\label{prop:covariance-closure}
Put $g_R=G_R-\bar G_R$. Suppose that
\begin{equation}\label{eq:covariance-tail-contract}
 \lim_{m\to\infty}\sup_R\sum_{k>m}
       \left\|\int g_R\otimes(g_R\circ(F_R^*)^k)\dd\nu_R^*\right\|=0.
\end{equation}
Suppose also that zero asymptotic variance of $v\cdot g_R$ implies a
coboundary whose transfer function has a representative for which the
coboundary identity telescopes on the five real-rank periods of
Corollary~\ref{cor:real-period-rank}. Then the Green--Kubo matrices
\[
 \Sigma_R=\int g_R\otimes g_R\dd\nu_R^*
  +\sum_{k\ge1}\int\bigl(g_R\otimes(g_R\circ F^k)
                    +(g_R\circ F^k)\otimes g_R\bigr)\dd\nu_R^*
\]
are continuous and satisfy $cI_4\le\Sigma_R\le CI_4$ on $I$ for some
$c,C>0$.
\end{proposition}
\begin{proof}
Corollary~\ref{cor:all-moment-continuity} makes every fixed-lag matrix
continuous, including the zero-lag term. The uniformly summable tail in
\eqref{eq:covariance-tail-contract} therefore gives continuity of the
series. The usual expansion of the covariance of $\sum_{j<n}g_R\circ F^j$
is a finite identity; dividing it by $n$ and using absolute summability
identifies its limit with $\Sigma_R$. Thus $\Sigma_R$ is positive
semidefinite and its directional quadratic form is the asymptotic
variance.

If $v^T\Sigma_Rv=0$, the stipulated coboundary identity gives on each of
the five selected periods
\[
 v_1K_1+v_2K_2+v_3N+v_4T-h(v\cdot\bar G_R)=0.
\]
The augmented vector $(v_1,v_2,v_3,-v\cdot\bar G_R,v_4)$ annihilates
their five-dimensional record matrix, whose determinant is nonzero by
Corollary~\ref{cor:real-period-rank}. Hence $v=0$. Each matrix is
positive definite. Continuity and compactness of $I$ give uniform lower
and upper eigenvalue bounds. The periodic representative premise is
used explicitly here; a merely measurable identity is insufficient.
\end{proof}

\begin{lemma}[An absolute raw-residual summation criterion]
\label{lem:raw-residual-sum}
Let $\mathcal W$ be countable. For $w\in\mathcal W$ assign a lattice
label $k_w\in\Z^3$ and a function $q_w\in L^1(\R)$ whose second
distributional derivative is a finite signed measure. Suppose
\[
 A_0=\sum_w\|q_w\|_1<\infty,\qquad
 A_2=\sum_w\|D^2q_w\|_{\TV}<\infty.
\]
Then
$H(u,b)=\sum_w e^{iu\cdot k_w}\widehat q_w(b)$ belongs to
$L^1(\T^3\times\R)$ and, for $B\ge1$,
\begin{align}\label{eq:absolute-residual-sum}
 \int_{\T^3\times\R}|H|&\le2(2\pi)^3(A_0+A_2),\\
 \int_{\T^3\times\{|b|>B\}}|H|&\le2(2\pi)^3A_2/B.\nonumber
\end{align}
The same estimates hold uniformly for a parameter family whenever the
corresponding sums are uniformly bounded.
\end{lemma}
\begin{proof}
Distributional differentiation and the stated Fourier convention give
$-b^2\widehat q_w(b)=\widehat{D^2q_w}(b)$. Thus
\[
 |\widehat q_w(b)|\le\min\{\|q_w\|_1,
                    |b|^{-2}\|D^2q_w\|_{\TV}\}\quad(b\ne0).
\]
Integrate the first estimate on $|b|\le1$ and the second on $|b|>1$,
and then sum by Tonelli. The torus has volume $(2\pi)^3$.
The series itself converges uniformly as a transform of the absolutely
summable $L^1$ densities. Integrating only $|b|>B$ gives the second
bound. All lattice multiplicities are included in the two sums; no
bound on the number of words has been suppressed.
\end{proof}

To use this lemma, the residuals must be those of a decomposition of the
actual full measure, including singular itinerary boundaries. A jump left
in $q_w$ contributes a derivative of a point mass to $D^2q_w$ and does
not satisfy the hypothesis. The individual jump estimates of
Proposition~\ref{prop:general-edge} do not prove the sum $A_2<\infty$.
Nor does the scalar tail \eqref{eq:exponential-return} bound inverse
Jacobian factors or the variation of a density near a critical value.
If $A_2$ grows with $n$, its actual growth must be used in the frequency
splice rather than treated as a fixed constant.

There is nevertheless a useful exact truncation budget. For an insertion
$|w|\le M_0$, deleting the physical words with $N_{n,R}>L$ changes its
record measure, and its raw transform in uniform norm, by at most
\begin{equation}\label{eq:physical-count-truncation}
 M_0C_1e^{-a_1L/n}.
\end{equation}
Indeed total variation contracts under pushforward, and
\eqref{eq:block-exponential-moment} bounds the discarded mass.
On a frequency band of volume $V$ the integrated error is at most
$VM_0C_1e^{-a_1L/n}$. Thus the sufficient physical-count budget for a
prescribed integrated error $\varepsilon$ is
\[
 L\ge\frac{n}{a_1}\max\{0,\log(VM_0C_1/\varepsilon)\}.
\]
This controls a finite-band error, not a pointwise density error or an
infinite frequency band. It is compatible with retaining the original
raw datum and with the separate all-branch derivative contract above.

\section{The joint local-limit conclusion and its realization}
\label{sec:final-realization}
The endpoint of this article remains the parameter-uniform raw mixed
density local limit \eqref{eq:LLT} for the physical record
$(K_{n,R},N_{n,R},T_{n,R})$, with the weighted positive-denominator
interface of Proposition~\ref{prop:conditioning}. Neither the return
section nor the recorded observables are replaced by independent flights,
a smoothed observation, or an unrelated geometric problem.

The change from $Y_R$ to $Y_R^*$ is explicit, and the periodic records
are those of its actual first return. The joint periodic annihilator is
trivial. The adjacent-increment comparison yields a uniform positive
period discrepancy with logarithmic collision length, and
Theorem~\ref{thm:one-step-coercivity} gives its positive-measure
consequence for controlled regular phases. The phase-reconstruction and
Fredholm hypotheses of Theorem~\ref{thm:phase-resolvent-criterion} state
precisely how these estimates would give an induced resolvent bound.
Their verification on the actual induced strong/weak spaces is still
required. No division by an uncontrolled spectral modulus is implicit.

The collision-space spectral input has a different, now completed use:
Theorem~\ref{thm:exponential-return} proves uniform exponential moments
of the genuine return block. It strengthens finite-record stability to
all finite moments. The common-chart construction in
Proposition~\ref{prop:finite-common-patches} gives an explicit
initial-state remainder for the moving-domain density comparison.
These facts are not substitutes for uniform summability of correlations;
the additional variance and cohomology requirements are exactly those
of Proposition~\ref{prop:covariance-closure}.

At the physical-clock level, the first-order functional law and its
geometric rate comparison remain valid. The stronger maximum estimate
in Theorem~\ref{thm:conditioned-maximal-block} now controls all unfinished
returns on a growing interval, at logarithmic scale, even after
conditioning of polynomially small positive probability. Corollary
\ref{cor:conditional-path-budget} supplies the corresponding explicit
square-root path error. Thus the unfinished-block premise no longer
needs to be assumed in that range of conditioning probabilities.
The functional Gaussian theorem for completed records, and the boundary
comparison when an exact lattice conditioning event is changed, are
separate inputs. This distinction is recorded in
Remark~\ref{rem:conditional-interface-scope}.

For the raw inversion itself, the exact edge subtraction and the local
edge criterion retain every correction term. The all-branch residual
must include central critical words and singular itinerary boundaries.
Lemma~\ref{lem:raw-residual-sum} specifies an absolute derivative-norm
criterion for that sum; no uncontrolled word multiplicity is hidden in
it. The required uniform sum has not been established here. Its growth
constants, together with the actual induced operator estimates, must
satisfy the strict splice inequalities \eqref{eq:splice} or directly
prove \eqref{eq:local-tail}. The weighted insertion class needs the same
verification. A positive conditioning denominator then follows from the
raw LLT on central windows by Proposition~\ref{prop:conditioning}.

The resulting dependency chain is exact: the geometry and periodic
coercivity, the genuine exponential return tail, the common finite-record
charts, and the conditioned unfinished-block estimate are proved;
induced phase reconstruction, the infinite covariance criterion, and
the global raw residual bounds are the remaining analytical work needed
to instantiate the full density theorem. No complete unconditional raw
LLT or full A3 density input is asserted before those premises are
verified. The proof of the new exponential return bound imports the
established collision spectral theory at its stated scope, not a theorem
about the unbounded induced observables.

\subsection*{Acknowledgment and reproducibility}
AI-assisted analysis contributed to this manuscript. The supplementary
source audit fixes the mathematical baseline, the referee report and
all imported source files; the response describes the revisions to each
numbered request. Finite computations check algebra and explicit orbit
margins and do not replace the analytical proofs. Independent human
specialist review is not claimed by this manuscript.

\clearpage
\begin{thebibliography}{9}
\bibitem{A2G} Q. Qi, \emph{Scalar collision laws and recognition of periodic dispersing billiards}, A2 v43, primary article and technical supplement, theta-theory repository, 2026.
\bibitem{SV} D. Sz\'asz and T. Varj\'u, \emph{Local limit theorem for the Lorentz process and its recurrence in the plane}, Ergodic Theory Dynam. Systems \textbf{24} (2004), 257--278.
\bibitem{DN} D. Dolgopyat and P. N\'andori, \emph{On mixing and the local central limit theorem for hyperbolic flows}, Ergodic Theory Dynam. Systems \textbf{40} (2020), 142--174; arXiv:1710.08568v1, Definition 2.2 and Proposition 6.3.
\bibitem{DZ} M. F. Demers and H.-K. Zhang, \emph{A functional analytic approach to perturbations of the Lorentz gas}, Comm. Math. Phys. \textbf{324} (2013), 767--830; arXiv:1210.1261.
\bibitem{BDL} V. Baladi, M. F. Demers and C. Liverani, \emph{Exponential decay of correlations for finite horizon Sinai billiard flows}, Invent. Math. \textbf{211} (2018), 39--177; arXiv:1506.02836v3, including the correction note.
\bibitem{Young} L.-S. Young, \emph{Statistical properties of dynamical systems with some hyperbolicity}, Ann. of Math. (2) \textbf{147} (1998), 585--650; author preprint ESI 445 (1997), Section 8, Theorem 6.
\bibitem{BHN} H. Bruin, M. Holland and M. Nicol, \emph{Liv\v{s}ic regularity for Markov systems}, Ergodic Theory Dynam. Systems \textbf{25} (2005), 1739--1765; doi:10.1017/S0143385705000179; arXiv:math/0503690.
\bibitem{DZ11} M. F. Demers and H.-K. Zhang, \emph{Spectral analysis of the transfer operator for the Lorentz gas}, J. Mod. Dyn. \textbf{5} (2011), 665--709; doi:10.3934/jmd.2011.5.665.
\end{thebibliography}

\end{document}
